\documentclass[11pt]{article}

\usepackage[utf8]{inputenc}
\usepackage[T1]{fontenc}
\usepackage{lmodern}
\usepackage{microtype}
\usepackage{geometry}
\geometry{margin=1in}
\setlength{\emergencystretch}{4em}
\usepackage{amsmath,amssymb,amsthm,mathtools}
\usepackage{booktabs,tabularx,array,longtable}
\usepackage{enumitem}
\usepackage{xcolor}
\usepackage{hyperref}
\usepackage[nameinlink,capitalize]{cleveref}

\hypersetup{
  colorlinks=true,
  linkcolor=black,
  citecolor=black,
  urlcolor=black,
  pdfauthor={E. S. Brooke},
  pdftitle={The Geometric Substrate as Cost Structure of Comparison Continuations}
}

\setlist[itemize]{leftmargin=1.4em, topsep=3pt, itemsep=2pt}
\setlist[enumerate]{leftmargin=1.6em, topsep=3pt, itemsep=2pt}
\renewcommand{\arraystretch}{1.15}

\newtheorem{theorem}{Theorem}[section]
\newtheorem{lemma}[theorem]{Lemma}
\newtheorem{proposition}[theorem]{Proposition}
\newtheorem{corollary}[theorem]{Corollary}
\newtheorem{definition}[theorem]{Definition}
\newtheorem{assumption}[theorem]{Assumption}
\newtheorem{principle}[theorem]{Principle}
\newtheorem{gate}[theorem]{Gate}
\theoremstyle{remark}
\newtheorem{remark}[theorem]{Remark}
\newtheorem{warning}[theorem]{Boundary}
\newtheorem{audit}[theorem]{Audit Note}

\newcommand{\G}{\Gamma}
\newcommand{\Ccap}{C}
\newcommand{\entails}{\Vdash}
\newcommand{\goes}{\rightsquigarrow}
\newcommand{\Cost}{\kappa}
\newcommand{\Cont}{\operatorname{Cont}}
\newcommand{\Erec}{E_{\rm rec}}
\newcommand{\phide}{\phi_{D\to E}}
\newcommand{\Req}{\phi_{\rm eq}}
\newcommand{\Kgeom}{\mathcal K}
\newcommand{\Mgeom}{\mathcal M}
\newcommand{\Dgeom}{\mathcal D}
\newcommand{\Egeom}{E_{\rm geom}}
\newcommand{\Jsrc}{\mathcal J}
\newcommand{\Status}[1]{\textsf{#1}}
\newcommand{\R}{\mathbb R}
\newcommand{\eps}{\varepsilon}
\newcommand{\dvol}{\,d^3x}
\newcommand{\Boxop}{\square}
\newcommand{\Lie}{\mathcal L}

\title{The Geometric Substrate as Cost Structure of Comparison Continuations\\
\large Finite-Continuability Translation of the Paper 6 GR Closure and the Weak-Field Trace Obstruction}
\author{E. S. Brooke}
\date{Version 1.4 -- Eternalist commitment for the metric-emergence ladder; structural-conditional reading of ``becomes available''}

\begin{document}
\maketitle

\begin{abstract}
This paper begins the geometric-substrate program for the Admissibility Physics Framework.  Paper 10 made continuation primitive: in context \(\G\), a distinction-expression \(D\) may continue as \(E\) within capacity bound \(C\), written \(\G\entails_C D\goes E\).  Geometry is then the regularized cost structure of comparison, localization, and transport continuations.  The present paper asks what must be proved for that slogan to become gravitational physics.

The answer is organized as a gate-based weak-field recovery program.  First, finite comparison continuations must admit a local cost germ.  Second, that germ must pass a quadraticity certificate so that a metric kernel is not smuggled.  Third, recruitment load from matter must perturb the clock and ruler sectors of that kernel.  Fourth, the perturbation must have the right scalar normal form.  Under the calibrated point-source load and proper-time dictionary, the clock sector conditionally recovers the Schwarzschild redshift factor; this is a real recovery, but it fixes only \(\Theta_t\), not the spatial response \(\Theta_s\).

The central technical result of this draft is the scalar weak-field closure package.  In each local clock-ruler comparison plane, the static scalar load response is represented by \(R\in GL^+(2,\mathbb R)\).  After coordinate gauge, anisotropic shear, and parity-forbidden time-space mixing are removed, the first-order scalar generator has two parameters.  Redshift calibrates the clock parameter.  The remaining trace is exactly the PPN-\(\gamma\) obstruction:
\[
        \theta=1-\gamma_C.
\]
Thus \(\gamma_C=1\) is equivalent to a traceless generator, a unimodular response, reciprocal clock-ruler scaling, and preservation of the local comparison cell.  A same-sign scalar depletion gives \(\gamma_C=-1\), showing why redshift recovery alone is insufficient.

The paper therefore reduces leading weak-field GR recovery to a precise next theorem target: derive the no-trace/unimodular scalar load-response gate from the geometric recruitment functional.  The full represented-geometry GR closure is not reproved here: it is imported from Paper~6 and its Supplement under their H1--H3/A9/Lovelock hypotheses.  What remains open in the present paper is the finite-continuability translation of those gates, especially the local no-trace scalar response that appears as the weak-field shadow of the Paper~6 closure.
\end{abstract}

\begin{center}
\fbox{\begin{minipage}{0.88\textwidth}
\centering
\vspace{3pt}
\textbf{Governing claim.}\\[3pt]
Geometry is not primitive background.  It is the represented cost structure of admissible comparison continuations.\\[4pt]
\[
\begin{aligned}
\G\entails_C x\goes y
&\Longrightarrow
\text{least-cost comparison profile}\;K_{x\to y}\\
&\Longrightarrow
\text{local quadratic cost kernel}\\
&\Longrightarrow
\text{metric representation, when the gates close.}
\end{aligned}
\]
\vspace{3pt}
\end{minipage}}
\end{center}

\tableofcontents

\section{Orientation: the geometric-substrate problem}

Paper 10 supplied a continuation-first language for APF.  Its central move was to separate admissible possibility from dynamics:
\[
        \G\entails_C D\goes E
\]
states that \(D\) can continue as \(E\) within capacity bound \(C\).  A dynamics, probability law, or variational rule is a later selector or weighting on admissible continuation words.

The same paper proposed a compressed reading of geometry:
\begin{quote}
Geometry is the regularized cost structure of carrying comparison, localization, and transport continuations.
\end{quote}
This paper begins the work required to make that sentence theorem-grade.

The geometric-substrate problem has five parts:
\begin{enumerate}
\item define comparison continuations operationally;
\item show when least-cost comparison behavior admits a smooth metric representation;
\item explain why the represented metric has Lorentzian rather than Euclidean signature;
\item derive how committed recruitment load changes the cost kernel for later comparison continuations;
\item align the weak-field cost-response calculation with the Paper~6 represented-geometry derivation of the Einstein equations.
\end{enumerate}

The present draft closes none of these globally.  It gives a clean starting architecture, proves the simplest representation lemmas under explicit hypotheses, and preserves the leading-order redshift calculation as the first worked recovery.

\begin{warning}[Status discipline]
Throughout this paper, \Status{P} means proved from stated definitions and assumptions in this paper; \Status{C} means conditional on an explicit gate or structural commitment; \Status{R} means recovered in a represented regime; \Status{I} means imported from another APF paper under its stated hypotheses; and \Status{O} means open.  The present document is mostly \Status{C/R/I}: it translates APF primitives into a geometric program, recovers the redshift sector under stated commitments, and imports the represented-geometry Einstein closure from Paper~6 rather than rederiving it.
\end{warning}

\subsection*{On operational language}

This paper develops geometry as the represented cost structure of finite comparison continuations.  The continuation primitive imported from Paper~10 is a static algebraic relation $\G\entails_C D\goes E$ on continuation profiles, and the metric-emergence ladder of \S\ref{sec:metric-emergence} derives the quadratic metric kernel as a representation \emph{conditional on a series of structural gates closing simultaneously} --- not as something the framework's primitives produce by temporal evolution.

The phrasings ``becomes available,'' ``emerges,'' ``closes,'' and ``recovers'' that appear throughout this paper are the local readings of conditional structural commitments.  Where the prose says ``the metric becomes available only after several gates close,'' the structural reading is ``the metric is a representation conditional on the gates' joint structural validity.''  Where the prose says ``admissible chart changes become diffeomorphism-type relabelings in the smooth represented regime,'' the structural reading is ``in the smooth represented regime, the admissible relabelings are exactly the diffeomorphism-type ones.''  The ``becomes'' is locator-language, not process-language.

Time itself is derived in this corpus: Paper~3 derives the arrow of time as the loss of global admissible coordination, and Paper~6 derives the metric and dimensionality of spacetime.  The geometric machinery developed here operates on the static comparison structure $(c_x, q_x)$ at the substrate; what reads as ``response,'' ``recruitment,'' or ``redshift recovery'' is the local reading of conditional structural commitments under temporal slicing.  A reader who wants the structural reading can take every ``becomes available,'' ``emerges,'' and ``recovers'' as shorthand for the local reading of the corresponding structural conditional.  The technical apparatus developed below --- the cost germ $c_x$, the quadratic kernel $q_x$, the geometric recruitment functional, the cost-curvature response --- is already structural; only the pedagogical wrapping is temporal.

This convention matches Paper~0's ``A note on language'' in the Prelude \S1 and Paper~0's Descriptive Reading chapter on temporal language and the eternalist commitment, Paper~1 supplement v8.24's ``On the use of operational language'' note in \S1, and Paper~10 v1.12's parallel note in \S1.  The framework is in this respect in the company of block-universe readings of general relativity, the Wheeler--DeWitt timeless quantum gravity programme, and the Page--Wootters mechanism for time emergence in entangled quantum systems --- structural at the foundation, with time as a derived feature of how the structure is read locally.

\section{Main results and non-claims of this draft}
\label{sec:main-results}

The document is organized around a deliberately narrow claim: finite comparison continuations can be represented by geometry only after several gates close, and the weak-field clock sector can be recovered before the full tensorial and nonlinear theory is known.  The following table states what is actually established in this draft and what remains only a target.

\begin{longtable}{p{0.24\textwidth}p{0.15\textwidth}p{0.50\textwidth}}
\toprule
Result & Status & Meaning \\
\midrule
Comparison-cost germ & \Status{C/P} & Given finite comparison continuations and refinement regularity, one obtains a local cost germ for infinitesimal comparisons. \\
Quadratic metric kernel & \Status{P} & If the local comparison cost satisfies the parallelogram certificate, polarization produces a quadratic kernel. \\
Point-source exterior profile & \Status{C/P} & An assumed isotropic quadratic exterior kernel gives the Green-function profile \(\chi(r)\sim1/r\).  The kernel itself is not yet derived. \\
Clock-sector redshift recovery & \Status{R} & With calibrated point-mass load, linear clock-stiffness response, and the proper-time dictionary, the Schwarzschild redshift factor is recovered. \\
Clock-only insufficiency & \Status{P/R} & Redshift and Newtonian acceleration constrain only \(\Theta_t\), not the spatial response \(\Theta_s\). \\
Scalar spatial-closure package & \Status{P/C} & In the gauge-reduced static scalar sector, the trace of the local clock-ruler response generator is exactly the PPN-\(\gamma\) obstruction. \\
Weak-field exterior recovery & \Status{C/R} & If the no-trace/unimodular scalar response gate closes, the first-post-Newtonian Schwarzschild exterior follows in isotropic weak-field form. \\
Full represented GR & \Status{I/C} & Imported from Paper~6 Supplement under H1--H3/A9 plus Lovelock uniqueness; finite-continuability translation of those gates remains an open Paper~9 task. \\
\bottomrule
\end{longtable}

\begin{audit}[Non-claims]
This paper does not rederive the Einstein equations, the value of \(G\), the dimension \((1+3)\), or the full equivalence principle.  Those represented-geometry closures are supplied by Paper~6 under its own hypotheses.  The present paper instead translates the same geometry into finite-comparison language and isolates the weak-field spatial coefficient as a single local scalar trace problem.  It states the exact gate whose closure would recover \(\gamma_{\rm PPN}=1\) without importing the full Paper~6 Lovelock route.
\end{audit}


\section{Alignment with Paper 6: represented-GR source of record}
\label{sec:paper6-alignment}

Paper~6 and its Technical Supplement already contain the represented-geometry route to GR.  In that paper the sequence is:
\[
\begin{aligned}
\text{cost distance}
&\Rightarrow g_{\mu\nu}
\Rightarrow (-,+,+,+),\ d=4 \\
&\Rightarrow \text{capacity-gradient curvature}
\Rightarrow \text{A9 closure} \\
&\Rightarrow \text{Lovelock uniqueness}
\Rightarrow
G_{\mu\nu}+\Lambda g_{\mu\nu}=\kappa T_{\mu\nu}.
\end{aligned}
\]
Paper~9 should therefore not duplicate that proof.  Its role is to express the same represented-geometry closure in the newer finite-continuability language and to identify what that closure looks like in the weak-field scalar sector.

\begin{theorem}[Paper 6 represented-GR import]
\label{thm:paper6-gr-import}
Under the Paper~6 hypotheses H1--H3, together with the A9 closure and Lovelock uniqueness in the represented smooth Lorentzian regime, the admissibility-preserving geometric dynamics are governed by
\[
        G_{\mu\nu}+\Lambda g_{\mu\nu}=\kappa T_{\mu\nu}.
\]
This theorem is imported here as a represented-geometry result, not reproved.
\end{theorem}

\begin{corollary}[Weak-field shadow of the Paper 6 closure]
\label{cor:paper6-weak-field-shadow}
If the Paper~6 represented-GR closure is accepted, then in the static weak-field exterior of a point mass the metric representation has the first-post-Newtonian Schwarzschild form
\[
        ds^2
        =-
        \left(1-\frac{2GM}{rc^2}\right)c^2dt^2
        +
        \left(1+\frac{2GM}{rc^2}\right)d\mathbf x^2
        +O(c^{-4}),
\]
and hence the scalar weak-field response of Paper~9 has
\[
        \gamma_C=1,
        \qquad
        \theta=1-\gamma_C=0.
\]
\end{corollary}

\begin{audit}[What Paper 9 adds beyond importing Paper 6]
The import theorem says that the represented smooth geometry satisfies the Einstein equations once the Paper~6 gates close.  Paper~9 adds a lower-level translation: redshift fixes the clock response, while the entire leading spatial ambiguity is the trace of a local clock-ruler scalar response generator.  Thus the no-trace/unimodular gate is the finite-continuability weak-field face of the Paper~6 A9/Lovelock closure.  If Paper~6 is used, \(\theta=0\) follows in the weak-field represented regime.  If Paper~9 is developed bottom-up, \(\theta=0\) is the next primitive substrate theorem to prove.
\end{audit}

\section{Reader path and theorem spine}
\label{sec:reader-path}

The draft has one central spine.  The early sections justify when a finite comparison calculus may be represented geometrically; the middle sections recover the clock sector; the weak-field section isolates the remaining spatial obstruction.

\begin{center}
\fbox{\begin{minipage}{0.9\textwidth}
\textbf{Publication spine.}
\[
\begin{aligned}
\text{finite comparison continuations}
&\Rightarrow \text{local cost germ} \\
&\Rightarrow \text{quadratic metric kernel, if parallelogram closes} \\
&\Rightarrow \Theta_t(r) \text{ from calibrated clock response} \\
&\Rightarrow \theta=1-\gamma_C \text{ in the scalar normal form} \\
&\Rightarrow \gamma_C=1 \text{ iff scalar response is traceless} \\
&\Rightarrow \text{conditional first-PN Schwarzschild exterior.}
\end{aligned}
\]
\end{minipage}}
\end{center}

The main mathematical reduction is therefore not ``derive GR'' in one jump.  It is the following finite list of closures:
\begin{enumerate}
\item prove the comparison-cost germ has a quadratic tangent limit;
\item prove the static point-source recruitment profile produces \(\Phi_C=-GM/r\) after calibration;
\item prove the clock response gives \(\Theta_t=1+2\Phi_C/c^2+O(c^{-4})\);
\item put the scalar clock-ruler response into gauge-reduced normal form;
\item prove the scalar trace vanishes: \(\operatorname{tr}L_{\rm sc}=0\).
\end{enumerate}
The first four items are represented or conditional in this draft.  The fifth is the new decisive target.

\begin{audit}[Why the trace problem is the right next target]
A redshift-only paper can always hide spatial curvature.  A metric-fitting paper can always insert \(\gamma=1\).  This paper does neither.  Once redshift fixes the clock coefficient, the scalar normal form proves that the entire first-order spatial ambiguity is the determinant/trace of the local clock-ruler response map.  Closing that trace gate is exactly what would turn the redshift calculation into a leading two-sector weak-field recovery.
\end{audit}


\section{Imports from finite continuability}

\begin{definition}[Continuation judgment]
The primitive APF judgment is
\[
        \G\entails_C D\goes E.
\]
It means that, in context \(\G\), distinction-expression \(D\) can continue as distinction-expression \(E\) using enforcement capacity no greater than \(C\).
\end{definition}

\begin{definition}[Continuation cost]
The cost of a continuation is
\[
        \Cost_\G(D\goes E)
        :=
        \inf\{C:\G\entails_C D\goes E\}.
\]
An identity-class continuation is a continuation with \(E\in[D]_\G\), where \([D]_\G\) is the tested continuation-equivalence class of \(D\).
\end{definition}

\begin{definition}[Recruitment profile]
For a continuation \(D\goes E\), the recruitment profile \(\phide(x)\) records where and how strongly substrate degrees of freedom are recruited to carry that continuation.  The associated recruitment cost is denoted
\[
        \Erec[\phi;\G].
\]
Fields are represented equilibrium recruitment profiles.  Radiation is cost discharge from failed adiabatic recruitment.  Geometry is the special case in which the recruited substrate carries comparison, localization, and transport continuations.
\end{definition}

\section{Comparison continuations}

\begin{definition}[Comparison continuation]
A comparison continuation from \(x\) to \(y\) in context \(\G\) is a substrate operation \(K:x\goes y\) that carries a probe distinction from \(x\) to \(y\), preserves the probe's identity class along the way, and establishes endpoint correspondence.  Its cost decomposes as
\[
\Cost_\G(K)
=
\Cost_{\rm transport}(K)
+
\Cost_{\rm identity}(K)
+
\Cost_{\rm corr}(K).
\]
The directed comparison cost is
\[
        d^+_\G(x,y)
        :=
        \inf\{\Cost_\G(K):K\text{ carries an admissible comparison continuation }x\goes y\}.
\]
\end{definition}

The decomposition is not meant to introduce three independent substances.  It is a bookkeeping separation: moving the probe, preserving what counts as the same probe, and matching endpoints are distinct burdens that may share substrate support.

\begin{definition}[Comparison substrate]
A comparison substrate \(\Mgeom\) is a regular domain of distinguishable comparison sites equipped with admissible comparison continuations between nearby sites.  It is \emph{locally comparison-complete} at scale \(\ell\) if for every pair \(x,y\in\Mgeom\) with represented separation below \(\ell\), at least one admissible comparison continuation connects \(x\) and \(y\).
\end{definition}

\begin{definition}[Local quadratic cost kernel]
A comparison substrate has a local quadratic cost kernel at \(x\) when infinitesimal comparison costs can be represented by a bilinear form \(q_x\) on tangent comparison directions:
\[
        \Cost_\G(x\goes x+dx)
        =
        \lambda_x\,\sqrt{|q_x(dx,dx)|}+o(\|dx\|)
\]
for some positive local normalization \(\lambda_x\).  Equivalently, the squared leading-order comparison cost is quadratic in \(dx\).
\end{definition}

\begin{gate}[Metric-representation gate]
A metric representation of comparison cost is available only if the following certificates close:
\begin{enumerate}
\item \textbf{local comparison completeness}: nearby points can be compared;
\item \textbf{smooth refinement}: comparison costs vary continuously under refinement of sites;
\item \textbf{quadratic tangent limit}: squared leading-order cost is represented by a nondegenerate bilinear form;
\item \textbf{path additivity}: finite comparison cost is approximated by integration of infinitesimal cost along admissible paths;
\item \textbf{signature certificate}: the bilinear form has one time-like and three space-like directions, or whatever signature is established by the substrate gates;
\item \textbf{dimension certificate}: the comparison tangent space has the represented dimension claimed.
\end{enumerate}
\end{gate}

\begin{proposition}[Metric representation under comparison certificates]
\label{prop:metric-rep}
If a comparison substrate is locally comparison-complete, admits a smooth nondegenerate local quadratic cost kernel \(q_x\), and satisfies path additivity, then its finite comparison costs admit a metric-type representation
\[
        L[\gamma]
        =
        \int_\gamma \sqrt{|q_{\gamma(s)}(\dot\gamma(s),\dot\gamma(s))|}\,ds
\]
up to local normalization.  Minimizing admissible comparison cost is represented by extremizing \(L[\gamma]\).
\end{proposition}

\begin{proof}
The tangent-limit hypothesis gives a nondegenerate quadratic form \(q_x\) for infinitesimal comparison increments.  Smooth refinement allows these local forms to be patched along admissible paths.  Path additivity states that the leading finite comparison cost is obtained as the refinement limit of sums of infinitesimal costs.  The refinement limit is precisely the displayed path functional.  Extremal comparison continuations are the represented least-cost paths.  No claim about Lorentzian signature follows until the signature certificate closes.
\end{proof}

\begin{remark}[What is proved here]
\Cref{prop:metric-rep} is a representation theorem under metric-like hypotheses.  It does not derive smoothness, dimension, or signature.  Its value is bookkeeping: it identifies exactly which gates must close before geometry can be used as a metric rather than a slogan.
\end{remark}

\section{Metric-emergence ladder: from cost germ to quadratic kernel}\label{sec:metric-emergence}

The phrase ``geometry is comparison cost'' hides several levels.  To avoid smuggling a metric, this paper distinguishes four stages:
\[
\begin{aligned}
\text{finite comparison preorder}
&\Longrightarrow \text{local cost germ}\\
&\Longrightarrow \text{Finsler-type length}\\
&\Longrightarrow \text{quadratic metric kernel.}
\end{aligned}
\]
Only the last stage licenses the usual tensor notation.

\begin{definition}[Local comparison cost germ]
Let \(x\in\Mgeom\).  A local comparison cost germ is the directional refinement limit
\[
        c_x(v):=
        \lim_{\epsilon\downarrow0}
        \frac{d^+_\G(x,x+\epsilon v)}{\epsilon},
\]
whenever the limit exists for a represented tangent direction \(v\).  The germ is \emph{regular} when it is finite on admissible directions, positive on nonzero spacelike comparison directions, and homogeneous under positive rescaling:
\[
        c_x(av)=a c_x(v),\qquad a>0.
\]
\end{definition}

\begin{gate}[Finsler-length gate]
A regular cost germ gives a Finsler-type comparison representation when each \(c_x\) is convex on the admissible cone and finite path costs are refinement limits of
\[
        L_c[\gamma]=\int c_{\gamma(s)}(\dot\gamma(s))\,ds.
\]
This gate is weaker than a metric gate: no quadratic form is yet implied.
\end{gate}

\begin{gate}[Quadraticity gate]
A Finsler-type comparison representation collapses to a quadratic metric representation only when the squared local germ satisfies the parallelogram identity on each admissible linearized comparison sector:
\[
        c_x(u+v)^2+c_x(u-v)^2
        =2c_x(u)^2+2c_x(v)^2.
\]
For Lorentzian sectors this identity is imposed separately on the linearized positive and negative norm sectors, with the null cone obtained as the boundary between them.
\end{gate}

\begin{proposition}[Quadratic kernel from the parallelogram certificate]
\label{prop:parallelogram-kernel}
Suppose a local comparison cost germ \(c_x\) is regular on a linearized comparison sector and satisfies the parallelogram identity.  Then the polarization formula
\[
        q_x(u,v)=\frac14\bigl(c_x(u+v)^2-c_x(u-v)^2\bigr)
\]
defines a symmetric bilinear form whose diagonal recovers the squared germ:
\[
        q_x(v,v)=c_x(v)^2.
\]
If \(q_x\) is nondegenerate and the signature gate closes, \(q_x\) is the local metric kernel of that sector.
\end{proposition}

\begin{proof}
The parallelogram identity is exactly the condition under which a squared norm is generated by an inner product through polarization.  Symmetry follows immediately from the polarization formula.  Additivity and homogeneity follow by the standard polarization argument applied to the regular homogeneous germ.  Nondegeneracy and signature are not consequences of the parallelogram identity; they are additional certificates.  With those certificates in place, the bilinear form is the local quadratic comparison kernel.
\end{proof}

\begin{audit}[No metric before the quadraticity gate]
A calculation that uses \(g_{ab}\), Christoffel symbols, curvature, geodesics, or the Einstein tensor has already passed into the represented metric regime.  Such calculations are legitimate bridge calculations only after the metric-representation and quadraticity gates have been named.  The APF derivation proper must either prove those gates or list them as inputs.
\end{audit}

\begin{definition}[Comparison atlas]
A comparison atlas is a family of operational coordinate charts built from mutually calibrated local comparison continuations.  Transition maps between charts are admissible when they preserve tested endpoint correspondences and do not change the least-cost comparison profiles except by allowed coordinate relabeling.  In the smooth represented regime, admissible chart changes become diffeomorphism-type relabelings.
\end{definition}

\section{Signature and dimension: what must be locked}

A Euclidean metric is not enough for physics.  Physical comparison has a causal structure: some continuations can be coordinated, some cannot, and irreversible records define an oriented order.  In APF language, Lorentzian signature is expected when the substrate has both a finite coordination speed and an irreversible record preorder.

\begin{principle}[Causal signature from finite coordination]
If comparison continuations have a finite maximum coordination throughput, then admissible comparison directions separate into causally reachable, null-boundary, and unreachable classes.  A Lorentzian representation is available when these classes are smoothly represented by one distinguished time-like direction and a cone of null comparison directions.
\end{principle}

\begin{principle}[Temporal orientation from irreversible record order]
The time orientation is not imposed by a coordinate label.  It is the represented direction of irreversible commitment: once a continuation profile has been coarsened into a record, admissible future continuations must preserve that committed record or pay irreversible overwrite cost.
\end{principle}

\begin{gate}[Signature gate]
The signature gate closes when the finite-coordination cone and the irreversible-record preorder can be represented by a smooth cone field with one distinguished time-like dimension and a nondegenerate transverse spatial sector.
\end{gate}

\begin{gate}[Dimension gate]
The dimension gate closes when the number of independent local comparison modes is fixed by substrate capacity, representation stability, and admissible interaction structure.  This draft does not derive the four-dimensional count.  It treats the \((1+3)\)-dimensional Lorentzian substrate as an input to the leading-order recovery.
\end{gate}

\section{The geometric recruitment functional}

The geometric substrate is not a passive manifold.  It is the substrate that carries comparison continuations.  Matter and stress-energy distinctions recruit that substrate; the recruitment load changes the future cost of comparison.

\begin{definition}[Effective geometric recruitment load]
Let \(\chi_{\rm rec}(x)\) denote the effective accumulated recruitment load carried by the geometric substrate at \(x\).  The notation \(\chi\), rather than \(\rho\), is deliberate: the leading exterior point-mass profile is potential-like, scaling as \(1/r\), not as a volumetric density.
\end{definition}

\begin{definition}[Quadratic geometric recruitment functional]
In the linear-response regime around an unloaded substrate, a scalar recruitment load field \(\chi\) sourced by a mass-energy distinction \(\Jsrc\) is represented by
\[
        \Erec[\chi]
        =
        \frac12\int_{\Mgeom} \Kgeom^{ij}\,\partial_i\chi\,\partial_j\chi\,d^3x
        -
        \int_{\Mgeom} \Jsrc\,\chi\,d^3x
        + O(\chi^3),
\]
where \(\Kgeom^{ij}\) is the unloaded spatial comparison stiffness.  The cubic and higher terms encode nonlinear self-coupling and are not fixed in this draft.
\end{definition}

\begin{assumption}[Isotropic exterior kernel]
For a static, isolated, spherically symmetric source in the exterior linear-response regime, the unloaded spatial comparison stiffness is isotropic:
\[
        \Kgeom^{ij}=K_s\delta^{ij}.
\]
\end{assumption}

\begin{proposition}[Exterior point-source profile in the linear-response regime]
\label{prop:point-source-profile}
Under the quadratic recruitment functional and isotropic exterior kernel, a point source \(\Jsrc(x)=\mu M\delta^{(3)}(x)\) has an exterior equilibrium profile
\[
        \chi_{\rm rec}(r)=\frac{\mu M}{4\pi K_s r}
\]
up to the normalization convention for \(\mu\).
\end{proposition}

\begin{proof}
Varying \(\Erec\) gives the Euler-Lagrange equation
\[
        -\partial_i(\Kgeom^{ij}\partial_j\chi)=\Jsrc.
\]
With \(\Kgeom^{ij}=K_s\delta^{ij}\), this becomes
\[
        -K_s\nabla^2\chi=\mu M\delta^{(3)}(x).
\]
Using \(\nabla^2(1/r)=-4\pi\delta^{(3)}(x)\), the exterior solution is
\[
        \chi_{\rm rec}(r)=\frac{\mu M}{4\pi K_s r},
\]
with any additive constant fixed by the boundary condition \(\chi\to0\) as \(r\to\infty\).
\end{proof}

\begin{remark}[What this does and does not close]
\Cref{prop:point-source-profile} closes only the Green's-function step once the quadratic isotropic kernel is assumed.  It does not derive the kernel from APF primitives, does not derive the source normalization \(\mu\), and does not derive Newton's constant.  It is nevertheless useful: it shows that the required \(1/r\) exterior profile is the natural equilibrium profile of the simplest geometric recruitment functional.
\end{remark}

\section{Cost-curvature response}

\begin{principle}[Linear cost-curvature response]
In the exterior weak-load regime, the time-comparison cost-curvature responds linearly to the effective accumulated geometric recruitment load:
\[
        K_t(x)=K_{t0}\bigl(1-\alpha\chi_{\rm rec}(x)\bigr),
\]
where \(K_{t0}\) is the unloaded time-comparison stiffness and \(\alpha\) is the coupling between committed load and residual time-comparison stiffness.
\end{principle}

The sign is physical: substrate already committed to carrying recruitment load has reduced residual stiffness for additional identity-carrying time comparisons.  In represented geometry, this reduced stiffness appears as slower local proper-time accumulation relative to the unloaded coordinate normalization.

\begin{warning}[Status of the response law]
The response law is a structural commitment in this draft.  A closed derivation would require a nonlinear geometric recruitment functional, expansion around the loaded equilibrium, and proof that the leading correction to \(K_{t0}\) is exactly \(-\alpha\chi_{\rm rec}\).  That is one of the main open tasks of this paper.
\end{warning}

\subsection{Minimal perturbative response model}

The following lemma does not close the response law from APF primitives.  It formalizes the exact kind of load-probe coupling that would make the law unavoidable in a weak-load effective theory.

\begin{definition}[Time-comparison probe]
Let \(u\) be a small normalized probe field representing an infinitesimal rest-frame identity-comparison tick.  In the unloaded substrate its leading quadratic cost is
\[
        E_t[u]=\frac12\int K_{t0}\,u^2\,d\lambda,
\]
where \(d\lambda\) is the local comparison-volume element appropriate to the clock sector.
\end{definition}

\begin{assumption}[Bilinear load-probe coupling]
In the weak-load regime, the leading coupling between accumulated geometric recruitment load \(\chi\) and the time-comparison probe \(u\) is
\[
        E_t[u;\chi]
        =\frac12\int \bigl(K_{t0}-\lambda\chi(x)\bigr)u(x)^2\,d\lambda
        +O(u^3,\chi^2u^2),
\]
with \(\lambda>0\).  The minus sign states residual-stiffness depletion: already committed geometric load reduces available stiffness for an additional identity-comparison tick.
\end{assumption}

\begin{proposition}[Linear response from bilinear load-probe coupling]
\label{prop:response-from-coupling}
Under the bilinear load-probe coupling, the local time-comparison stiffness at fixed load is
\[
        K_t(x)=K_{t0}-\lambda\chi(x)+O(\chi^2),
\]
or equivalently
\[
        K_t(x)=K_{t0}\bigl(1-\alpha\chi(x)\bigr)+O(\chi^2),
        \qquad \alpha:=\lambda/K_{t0}.
\]
\end{proposition}

\begin{proof}
The local stiffness is the second variation of the probe cost with respect to the probe amplitude at the loaded equilibrium:
\[
        K_t(x)=\left.\frac{\delta^2 E_t[u;\chi]}{\delta u(x)^2}\right|_{u=0}.
\]
Taking the second variation of the displayed effective cost gives \(K_{t0}-\lambda\chi(x)\) to first order in \(\chi\).  Dividing by \(K_{t0}\) gives the normalized coefficient \(\alpha=\lambda/K_{t0}\).
\end{proof}

\begin{audit}[What the response lemma buys]
\Cref{prop:response-from-coupling} converts the response law from a naked ansatz into a precise perturbative target.  The remaining open work is now sharper: derive the negative bilinear load-probe coupling from the nonlinear geometric substrate functional, rather than merely assert the final linear response.
\end{audit}

\begin{definition}[Newtonian calibration]
In physical units, define the exterior point-mass calibration by
\[
        \alpha\chi_{\rm rec}(r)=\frac{2GM}{rc^2}.
\]
Equivalently, the combination \(\alpha\mu/(4\pi K_s)\) is fixed to \(2G/c^2\).  In this draft, \(G\) is a substrate-specific dimensional input, not an internally derived number.
\end{definition}

\begin{corollary}[Schwarzschild time factor as calibrated cost curvature]
\label{cor:schwarzschild-factor}
Under the point-source profile and Newtonian calibration,
\[
        K_t(r)=K_{t0}\left(1-\frac{2GM}{rc^2}\right)
\]
in the exterior linear-response regime.
\end{corollary}

\section{Proper time as normalized identity-continuation cost}

\begin{definition}[Proper-time dictionary]
In a represented geometric regime, the proper-time element for a static identity-carrying continuation is the normalized cost length of that continuation:
\[
        d\tau(r)=\sqrt{K_t(r)/K_{t0}}\,dt.
\]
This dictionary treats \(dt\) as the unloaded coordinate comparison parameter and \(d\tau\) as the local identity-continuation clock tick.
\end{definition}

\begin{remark}[Why the square root]
The local line element is quadratic in infinitesimal comparison displacement.  Cost-curvature multiplies the quadratic form; the represented cost length is therefore the square root of the curvature-normalized quadratic element.  This is a representation dictionary, not yet a derivation from a fully closed substrate action.
\end{remark}

\begin{definition}[Clock sector]
The \emph{clock sector} of a represented geometry is the part of the comparison-cost kernel that controls static rest-frame identity continuations.  In a static coordinate representation it is encoded by a lapse factor
\[
        N(x):=\frac{d\tau(x)}{dt}.
\]
The redshift calculation constrains this sector only.
\end{definition}

\begin{definition}[Spatial comparison sector]
The \emph{spatial comparison sector} is the part of the comparison-cost kernel that controls endpoint comparison at fixed clock label.  In a weak-field isotropic representation one may write, schematically,
\[
        ds^2=-N(x)^2c^2dt^2 + S(x)^2\,d\mathbf x^2,
\]
where redshift fixes the lapse factor \(N\) but not the spatial scale factor \(S\).
\end{definition}

\begin{proposition}[Redshift does not fix full geometry]
\label{prop:redshift-not-full-geometry}
Any two static represented geometries with the same lapse factor \(N(r)=d\tau(r)/dt\) give the same static gravitational redshift between radii, even if their spatial comparison sectors differ.  Therefore a redshift recovery alone cannot establish full GR.
\end{proposition}

\begin{proof}
For static emitter and receiver, local frequency is inverse local proper period.  The same coordinate period is compared against the two local lapse factors, so
\[
        \frac{\omega_A}{\omega_B}=\frac{N(r_B)}{N(r_A)}.
\]
No spatial comparison factor enters this expression.  Hence all represented metrics with the same lapse produce the same redshift.
\end{proof}

\begin{proposition}[Leading-order redshift recovery from cost-curvature response]
\label{prop:redshift-recovery}
Consider a static, spherically symmetric recruitment load generated by a point mass \(M\), in the exterior linear-response regime.  Suppose
\[
        K_t(r)=K_{t0}\bigl(1-\alpha\chi_{\rm rec}(r)\bigr),
        \qquad
        \alpha\chi_{\rm rec}(r)=\frac{2GM}{rc^2},
\]
and suppose proper time is represented by normalized identity-continuation cost:
\[
        d\tau(r)=\sqrt{K_t(r)/K_{t0}}\,dt.
\]
Then
\[
        d\tau(r)=\sqrt{1-\frac{2GM}{rc^2}}\,dt.
\]
Consequently, for a photon emitted by a static observer at \(r_B\) and received by a static observer at \(r_A>r_B\),
\[
        \frac{\omega_A}{\omega_B}
        =
        \sqrt{\frac{1-2GM/r_Bc^2}{1-2GM/r_Ac^2}}.
\]
In the weak-field limit,
\[
        \frac{\omega_A}{\omega_B}
        \approx
        1-
        \frac{GM}{c^2}\left(\frac1{r_B}-\frac1{r_A}\right).
\]
\end{proposition}

\begin{proof}
The first displayed formula follows immediately from substituting the calibrated recruitment load into the cost-curvature response law and then applying the proper-time dictionary.  For a static emitter and receiver, the same coordinate-period label is compared against two different local proper-time rates.  Local frequency is inverse local proper period, hence
\[
        \omega_A=\omega_B\frac{d\tau(r_B)}{d\tau(r_A)}.
\]
Substitution gives the redshift ratio.  Expanding \(\sqrt{1-2GM/rc^2}\) to first order gives the weak-field expression.
\end{proof}

\begin{warning}[Boundary of the result]
\Cref{prop:redshift-recovery} is a conditional recovery, not a closed derivation of GR.  It assumes: (i) the linear cost-curvature response, (ii) the exterior point-source recruitment profile, (iii) the proper-time-as-normalized-cost dictionary, and (iv) \(G\) as the substrate coupling in physical units.  Deriving these from substrate primitives, extending the result beyond the exterior linear-response regime, and recovering the full Einstein equations are open.
\end{warning}

\section{No-smuggling audit for the redshift recovery}

The redshift recovery is useful only if its inputs are named without disguise.  The calculation has four layers.

\begin{longtable}{p{0.23\textwidth}p{0.16\textwidth}p{0.50\textwidth}}
\toprule
Layer & Status & Content \\
\midrule
Continuation grammar & \Status{Imported} & The judgment \(\G\entails_C D\goes E\) and the interpretation of geometry as comparison-continuation cost are imported from Paper 22. \\
Exterior profile & \Status{Cond./proved} & The \(1/r\) profile is proved here only after assuming an isotropic quadratic exterior kernel. \\
Clock response & \Status{Assumed} & \(K_t=K_{t0}(1-\alpha\chi)\) is a linear-response commitment. Its sign and linearity are physically motivated but not derived from a nonlinear substrate action. \\
Unit calibration & \Status{Input} & The equality \(\alpha\chi=2GM/rc^2\) imports the physical coupling \(G\). The calculation recovers the form of the redshift, not the numerical value of \(G\). \\
Proper-time dictionary & \Status{Representational} & \(d\tau=\sqrt{K_t/K_{t0}}dt\) is the metric-regime dictionary between normalized identity-continuation cost and clock time. \\
\bottomrule
\end{longtable}

\begin{audit}[What is actually earned]
The calculation earns the following claim: if a point mass produces a potential-like recruitment load, if time-comparison stiffness is reduced linearly by that load, and if proper time is normalized cost length, then the static clock-redshift sector has the Schwarzschild form at leading exterior order.  It does not earn the spatial metric, the equivalence principle, nonlinear field equations, or a derived value of \(G\).
\end{audit}

\begin{audit}[What would close the result]
To turn the recovery into a derivation, one must derive the isotropic Green's-function kernel, the universal response law, the proper-time dictionary, and the physical coupling normalization from APF substrate primitives.  These are not cosmetic gaps; they are the main mathematical work of Paper 9.
\end{audit}

\section{Newtonian acceleration as a cost-gradient effect}

The redshift calculation recovers the clock sector.  The next target is motion.  In represented weak-field geometry, slow test bodies accelerate along gradients of the same potential that controls \(g_{tt}\).  Cost-theoretically, this should emerge because least-cost identity-carrying worldlines bend toward lower continuation cost.

\begin{definition}[Effective cost potential]
Define the effective cost potential \(\Phi_C\) by
\[
        \frac{K_t(r)}{K_{t0}}=1+\frac{2\Phi_C(r)}{c^2}.
\]
For the calibrated point-source profile,
\[
        \Phi_C(r)=-\frac{GM}{r}.
\]
\end{definition}

\begin{proposition}[Newtonian acceleration under the geodesic-cost dictionary]
Assume the metric-representation gate closes and slow test-body worldlines extremize normalized identity-continuation cost.  In the weak-field, slow-velocity regime with
\[
        \frac{K_t}{K_{t0}}=1+\frac{2\Phi_C}{c^2},
\]
the represented coordinate acceleration is
\[
        \mathbf a=-\nabla\Phi_C.
\]
For \(\Phi_C=-GM/r\), this gives the inverse-square law.
\end{proposition}

\begin{proof}[Proof sketch]
Under the geodesic-cost dictionary, the action for a slow test body is proportional to
\[
        \int d\tau
        \approx
        \int \left(1+\frac{\Phi_C}{c^2}-\frac{v^2}{2c^2}\right)dt
\]
up to an irrelevant constant normalization and sign convention.  Extremizing is equivalent to extremizing the ordinary weak-field Lagrangian \(L=v^2/2-\Phi_C\), yielding \(\mathbf a=-\nabla\Phi_C\).  The step is conditional on the metric and geodesic-cost dictionaries.
\end{proof}

\begin{remark}[Why this matters]
The same scalar cost load controls both clock rate and slow-body acceleration.  This is the beginning of an equivalence-principle-like universality claim, but it is not yet a derivation of the equivalence principle.  To close that claim, the response law must be shown universal for all admissible rest-frame identity continuations, independent of matter composition.
\end{remark}


\section{Weak-field two-sector kernel and the missing spatial response}

The clock-sector recovery should be tested against the standard weak-field hierarchy.  The redshift calculation gives a lapse.  A gravitational theory also needs the spatial comparison response.  The next formal object is therefore not a scalar clock stiffness alone, but a two-sector normalized cost kernel.

\begin{definition}[Normalized two-sector weak-field cost kernel]
\label{def:two-sector-kernel}
In a static, isotropic, weak-field represented regime, write the normalized comparison-cost line element as
\[
        ds_C^2
        =-\Theta_t(r)c^2dt^2+\Theta_s(r)d\mathbf x^2+O(c^{-4}),
\]
where
\[
        \Theta_t(r):=\frac{K_t(r)}{K_{t0}},
        \qquad
        \Theta_s(r):=\frac{K_s(r)}{K_{s0}}.
\]
The redshift result fixes \(\Theta_t\).  It does not fix \(\Theta_s\).
\end{definition}

\begin{definition}[Weak-field cost potential]
For a static exterior point source, define the cost potential \(\Phi_C\) by
\[
        \Theta_t(r)=1+\frac{2\Phi_C(r)}{c^2}+O(c^{-4}).
\]
Under the Newtonian calibration of the scalar recovery,
\[
        \Phi_C(r)=-\frac{GM}{r}.
\]
\end{definition}

\begin{definition}[Spatial response parameter]
The weak-field spatial comparison response is defined by
\[
        \Theta_s(r)=1-2\gamma_C\frac{\Phi_C(r)}{c^2}+O(c^{-4}).
\]
The dimensionless coefficient \(\gamma_C\) is the APF/cost-theoretic counterpart of the standard PPN parameter \(\gamma_{\rm PPN}\).  General relativity has \(\gamma_C=1\).
\end{definition}

\begin{proposition}[Clock observables fix only the lapse]
\label{prop:clock-only-lapse}
Gravitational time dilation, static redshift, and Newtonian slow-body acceleration determine the leading clock potential \(\Phi_C\) but do not determine the spatial-response coefficient \(\gamma_C\).
\end{proposition}

\begin{proof}
The redshift and time-dilation formulae depend only on the lapse factor \(\Theta_t^{1/2}\).  The Newtonian slow-motion limit follows from the same lapse through the weak-field particle action.  The spatial factor \(\Theta_s\) enters null propagation and spatial-curvature observables at the same post-Newtonian order, but it does not enter the static redshift ratio and does not affect the lowest-order slow-velocity acceleration law.  Therefore the clock sector fixes \(\Phi_C\) and leaves \(\gamma_C\) free.
\end{proof}

\begin{gate}[Spatial-response / PPN-\(\gamma\) gate]
\label{gate:spatial-response}
To recover the full leading weak-field Schwarzschild sector, the geometric substrate must determine the spatial comparison response
\[
        \Theta_s(r)=1-2\frac{\Phi_C(r)}{c^2}+O(c^{-4}),
\]
that is,
\[
        \gamma_C=\gamma_{\rm PPN}=1.
\]
Until this gate closes, the paper has recovered the clock sector and Newtonian limit, but not light bending, Shapiro delay, or the full first-post-Newtonian exterior metric.

\end{gate}

\begin{gate}[Reciprocal clock-ruler response]
\label{gate:reciprocal-clock-ruler}
The spatial-response gate can be sharpened into a candidate closure principle.  In the exterior weak-field scalar-load regime, the time-comparison and spatial-comparison sectors are required to be reciprocal to first post-Newtonian order:
\[
        \Theta_s(r)\Theta_t(r)=1+O(c^{-4}),
\]
or equivalently
\[
        \Theta_s(r)=\Theta_t(r)^{-1}+O(c^{-4}).
\]
Physical reading: clock comparison and ruler comparison are conjugate faces of the same finite-coordination substrate.  A load that reduces the residual rate of identity-carrying clock continuation must increase the cost of spatial endpoint comparison by the inverse factor if the same local null coordination boundary is to be used by all admissible probes.  This is a gate, not a theorem: it must ultimately be derived from the microscopic comparison functional.
\end{gate}

\begin{proposition}[Reciprocal response implies the PPN spatial coefficient]
\label{prop:reciprocal-implies-gamma}
Assume the two-sector weak-field form
\[
        \Theta_t(r)=1+\frac{2\Phi_C(r)}{c^2}+O(c^{-4}),
        \qquad
        \Theta_s(r)=1-2\gamma_C\frac{\Phi_C(r)}{c^2}+O(c^{-4}).
\]
If the reciprocal clock-ruler gate holds, then
\[
        \gamma_C=1.
\]
\end{proposition}

\begin{proof}
Multiplying the two sector factors gives
\[
\begin{aligned}
        \Theta_s\Theta_t
        &=\left(1-2\gamma_C\frac{\Phi_C}{c^2}\right)
          \left(1+2\frac{\Phi_C}{c^2}\right)+O(c^{-4})\\
        &=1+2(1-\gamma_C)\frac{\Phi_C}{c^2}+O(c^{-4}).
\end{aligned}
\]
The reciprocal gate requires this product to be unity through order \(c^{-2}\).  For a nonzero exterior potential this forces \(1-\gamma_C=0\), hence \(\gamma_C=1\).
\end{proof}

\begin{audit}[Why naive same-sign scalar depletion fails]
A tempting first guess is that every comparison stiffness is depleted in the same scalar way:
\[
        \Theta_s(r)=\Theta_t(r)=1+\frac{2\Phi_C(r)}{c^2}+O(c^{-4}).
\]
Compared with the PPN form \(\Theta_s=1-2\gamma_C\Phi_C/c^2\), this gives \(\gamma_C=-1\).  Thus redshift would still be recovered from the clock sector, but the spatial-curvature contribution needed for the standard two-sector weak-field observables would have the wrong coefficient.  The geometric substrate cannot be modeled as uniform same-sign depletion of every comparison channel.  The spatial sector must be conjugate, tensorial, or otherwise constrained; this is precisely the content of \Cref{gate:reciprocal-clock-ruler} and \Cref{gate:spatial-response}.
\end{audit}

\subsection{Candidate derivation route: clock-ruler comparison-cell conservation}

The reciprocal clock-ruler condition should not remain a slogan.  The following definitions isolate a sharper operational certificate.  The certificate is still a gate, but it is weaker and more testable than simply assuming \(\gamma_C=1\): it says that, direction by direction, the elementary comparison cell built from one clock tick and one spatial endpoint ruler is preserved by a static scalar load to first post-Newtonian order.

\begin{definition}[Directionwise clock-ruler comparison cell]
Fix a represented spatial comparison direction \(e\) in the static weak-field regime.  The normalized clock-ruler comparison cell in that direction is
\[
        \mathcal A_e(r)
        :=
        \left(\frac{d\tau(r)}{dt}\right)
        \left(\frac{d\ell_e(r)}{dx_e}\right)
        =
        \sqrt{\Theta_t(r)\Theta_{s,e}(r)},
\]
where \(d\ell_e=\sqrt{\Theta_{s,e}}\,dx_e\) is the represented spatial comparison length in direction \(e\).  In the isotropic weak-field sector, \(\Theta_{s,e}=\Theta_s\) for all spatial directions.
\end{definition}

\begin{gate}[Directionwise comparison-cell conservation]
\label{gate:cell-conservation}
In an exterior, static, scalar-load regime with no shear, torsion, anisotropic substrate recruitment, or composition-dependent comparison channel, the elementary clock-ruler comparison cell is conserved to first post-Newtonian order:
\[
        \mathcal A_e(r)=1+O(c^{-4})
\]
for every represented spatial direction \(e\).

Physical reading: a scalar geometric load may redistribute comparison capacity between clock-rate continuation and endpoint-ruler continuation, but it does not create or destroy the local directionwise comparison cell.  Clock slowing and ruler stretching are conjugate responses of the same finite-coordination substrate.
\end{gate}

\begin{proposition}[Cell conservation implies reciprocal clock-ruler response]
\label{prop:cell-implies-reciprocal}
In the static isotropic two-sector weak-field kernel, directionwise comparison-cell conservation implies
\[
        \Theta_s(r)\Theta_t(r)=1+O(c^{-4}).
\]
Consequently, under the PPN expansion of \Cref{prop:reciprocal-implies-gamma}, it implies \(\gamma_C=1\).
\end{proposition}

\begin{proof}
In the isotropic sector, \(\mathcal A_e=\sqrt{\Theta_t\Theta_s}\) is independent of \(e\).  The cell-conservation gate gives \(\sqrt{\Theta_t\Theta_s}=1+O(c^{-4})\).  Squaring gives \(\Theta_t\Theta_s=1+O(c^{-4})\).  The implication \(\gamma_C=1\) then follows from \Cref{prop:reciprocal-implies-gamma}.
\end{proof}

\begin{audit}[What the cell-conservation route does not prove]
Cell conservation is not the local speed of light, coordinate invariance, or redshift in disguise.  Local null measurements give a universal measured speed once the metric representation is assumed; they do not by themselves fix \(\Theta_s\).  The new content is the claim that a static scalar recruitment load preserves the directionwise clock-ruler comparison cell at order \(c^{-2}\).  This can fail if the load creates anisotropic shear, if clock and ruler comparisons draw on independent substrate reservoirs, if matter species couple to different comparison kernels, or if the geometric load changes the local comparison-cell budget rather than redistributing it.  These are empirical and mathematical failure modes, not notational details.
\end{audit}


\subsection{Sharper candidate route: canonical comparison pairing}

The cell-conservation gate can itself be decomposed.  A comparison between events has two faces: a primal clock-ruler displacement and a dual frequency-wavevector probe.  The dimensionless action/phase pairing between these faces is the invariant bookkeeping object.  This subsection isolates the extra condition needed to make that pairing imply reciprocal clock-ruler response.

\begin{definition}[Local canonical comparison pair]
Fix a spatial comparison direction \(e\) in a represented local frame.  The primal comparison variables are
\[
        q_e=(\tau,\ell_e),
\]
where \(\tau\) is clock continuation and \(\ell_e\) is endpoint-ruler continuation in direction \(e\).  The dual probe variables are
\[
        p_e=(\omega,k_e),
\]
where \(\omega\) is frequency cost and \(k_e\) is directional wave-number cost.  The local comparison one-form is
\[
        \vartheta_e=-\omega\,d\tau+k_e\,d\ell_e .
\]
Equivalently, \(\vartheta_e/\hbar\) is the local dimensionless phase increment of a probe used to test the comparison.
\end{definition}

\begin{gate}[Canonical load transport]
\label{gate:canonical-load-transport}
A static scalar recruitment load acts canonically on the local comparison pair if, direction by direction, its transport map preserves the comparison one-form up to terms of order \(c^{-4}\):
\[
        T_\chi^*\vartheta_e=\vartheta_e+O(c^{-4}).
\]
For a diagonal weak-field response this means that if
\[
        d\tau\mapsto a_t\,d\tau,
        \qquad
        d\ell_e\mapsto a_s\,d\ell_e,
\]
then the dual variables transform inversely,
\[
        \omega\mapsto a_t^{-1}\omega,
        \qquad
        k_e\mapsto a_s^{-1}k_e,
\]
so that the phase/action pairing is not changed by relabeling the same comparison through the loaded substrate.
\end{gate}

\begin{audit}[Canonical pairing alone is not enough]
\label{audit:canonical-not-enough}
\Cref{gate:canonical-load-transport} preserves the action/phase pairing, but it does not by itself force \(a_ta_s=1\).  Any nonzero product \(a_ta_s\) can be paired with inverse dual scaling and still preserve \(\vartheta_e\).  Therefore the spatial coefficient \(\gamma_C=1\) requires one additional physical statement: the scalar load may exchange clock and ruler comparison capacity, but it may not dilate the elementary primal clock-ruler comparison cell at order \(c^{-2}\).
\end{audit}

\begin{gate}[Trace-free clock-ruler load transport]
\label{gate:trace-free-transport}
In the exterior static scalar-load regime, after removing coordinate relabelings and anisotropic shear, the primal action of the load on each local clock-ruler comparison plane is trace-free to first post-Newtonian order.  Equivalently, there exists a scalar load parameter \(\sigma(r)=O(c^{-2})\) such that
\[
        d\tau\mapsto e^{\sigma(r)}d\tau,
        \qquad
        d\ell_e\mapsto e^{-\sigma(r)}d\ell_e
\]
for every spatial direction \(e\), up to \(O(c^{-4})\).

Physical reading: a static scalar geometric load does not create an independent local comparison-cell resource.  It transfers comparison capacity between conjugate clock and ruler channels.  The trace part would represent creation or destruction of the directionwise comparison cell; the trace-free part represents redistribution.
\end{gate}

\begin{proposition}[Trace-free transport implies cell conservation]
\label{prop:trace-free-implies-cell}
Assume the represented two-sector weak-field kernel and trace-free clock-ruler load transport.  Then directionwise clock-ruler comparison-cell conservation holds:
\[
        \mathcal A_e(r)=1+O(c^{-4}).
\]
Consequently,
\[
        \Theta_s(r)\Theta_t(r)=1+O(c^{-4})
        \quad\text{and}\quad
        \gamma_C=1.
\]
\end{proposition}

\begin{proof}
By definition, the normalized clock factor is \(d\tau/dt=\sqrt{\Theta_t}\), and the normalized ruler factor in direction \(e\) is \(d\ell_e/dx_e=\sqrt{\Theta_{s,e}}\).  Trace-free transport gives these factors as \(e^{\sigma}\) and \(e^{-\sigma}\), up to \(O(c^{-4})\).  Hence
\[
        \mathcal A_e
        =\left(\frac{d\tau}{dt}\right)
         \left(\frac{d\ell_e}{dx_e}\right)
        =e^{\sigma}e^{-\sigma}+O(c^{-4})
        =1+O(c^{-4}).
\]
The remaining implications are exactly \Cref{prop:cell-implies-reciprocal,prop:reciprocal-implies-gamma}.
\end{proof}

\begin{corollary}[Nested spatial-gate chain]
\label{cor:nested-spatial-chain}
The leading spatial coefficient can be closed by any of the following increasingly primitive certificates:
\[
\begin{aligned}
        \text{trace-free clock-ruler load transport}
        &\Longrightarrow
        \text{directionwise cell conservation}\\
        &\Longrightarrow
        \Theta_s\Theta_t=1+O(c^{-4})\\
        &\Longrightarrow
        \gamma_C=1.
\end{aligned}
\]
Thus the next mathematical target is sharper than \(\gamma_C=1\): derive trace-free transport from the geometric recruitment functional, or exhibit a countermodel showing that APF permits a trace component.
\end{corollary}


\subsection{Variational route: fixed-cell response and the trace obstruction}

Trace-free transport should itself be generated by a local response principle, not merely imposed.  The cleanest way to state the target is to decompose the diagonal clock-ruler load response into its trace and trace-free parts.

\begin{definition}[Logarithmic clock-ruler response variables]
\label{def:log-response-variables}
Fix a represented spatial direction \(e\).  Write the local diagonal load response as
\[
        \frac{d\tau}{dt}=a_t=e^{\xi_t},
        \qquad
        \frac{d\ell_e}{dx_e}=a_s=e^{\xi_s}.
\]
Define the trace and exchange variables
\[
        u_e:=\xi_t+\xi_s,
        \qquad
        v_e:=\xi_t-\xi_s.
\]
Then \(u_e=\log\mathcal A_e\) is the logarithmic clock-ruler cell dilation, while \(v_e\) is the trace-free exchange between clock and ruler comparison channels.
\end{definition}

\begin{gate}[Fixed-cell variational response]
\label{gate:fixed-cell-variational-response}
In the exterior static scalar-load regime, after removing coordinate relabelings and shear, the local response of a directionwise clock-ruler comparison plane is governed to quadratic order by an effective potential of the form
\[
        \mathcal V_e(u_e,v_e;\chi)
        =
        \frac12 A_e u_e^2
        +\frac12 B_e\bigl(v_e-\beta_e\chi\bigr)^2
        +O\bigl(|u_e|^3+|v_e|^3+\chi^2|u_e|+\chi^2|v_e|\bigr),
\]
with \(A_e,B_e>0\).  The scalar load sources the exchange mode \(v_e\) but does not source the trace mode \(u_e\) at first order.

Physical reading: committed geometric recruitment may shift capacity between clock continuation and ruler comparison, but a scalar static load has no independent channel for creating or destroying the elementary clock-ruler comparison cell.  Any first-order source term proportional to \(\chi u_e\) would be a new local comparison-resource reservoir and must be separately justified.
\end{gate}

\begin{proposition}[Fixed-cell variational response implies trace-free transport]
\label{prop:fixed-cell-implies-trace-free}
Assume the fixed-cell variational response gate and \(\chi=O(c^{-2})\).  The local equilibrium response satisfies
\[
        u_e=O(\chi^2)=O(c^{-4}),
        \qquad
        v_e=\beta_e\chi+O(\chi^2).
\]
Consequently
\[
        \mathcal A_e=e^{u_e}=1+O(c^{-4}),
\]
so directionwise comparison-cell conservation holds.  In the isotropic two-sector weak-field regime this implies
\[
        \Theta_s\Theta_t=1+O(c^{-4}),
        \qquad
        \gamma_C=1.
\]
\end{proposition}

\begin{proof}
Equilibrium requires \(\partial\mathcal V_e/\partial u_e=0\) and \(\partial\mathcal V_e/\partial v_e=0\).  To first order,
\[
        \frac{\partial\mathcal V_e}{\partial u_e}=A_eu_e+O(\chi^2),
        \qquad
        \frac{\partial\mathcal V_e}{\partial v_e}=B_e(v_e-\beta_e\chi)+O(\chi^2).
\]
Since \(A_e,B_e>0\), the equilibrium solution is
\[
        u_e=O(\chi^2),
        \qquad
        v_e=\beta_e\chi+O(\chi^2).
\]
Because \(\chi=O(c^{-2})\), the trace is absent through order \(c^{-2}\).  Hence \(\mathcal A_e=e^{u_e}=1+O(c^{-4})\).  The final implications are \Cref{prop:cell-implies-reciprocal,prop:reciprocal-implies-gamma}.
\end{proof}

\begin{definition}[Trace-obstruction coefficient]
\label{def:trace-obstruction-coefficient}
If the local response potential contains a first-order trace source,
\[
        \mathcal V_e(u_e,v_e;\chi)
        =
        \frac12 A_eu_e^2+\frac12 B_e(v_e-\beta_e\chi)^2
        -\lambda_e\chi u_e+O(3),
\]
then the equilibrium trace is
\[
        u_e=\eta_e\chi+O(\chi^2),
        \qquad
        \eta_e:=\lambda_e/A_e.
\]
The coefficient \(\eta_e\) is the first-order trace obstruction.
\end{definition}

\begin{proposition}[Trace obstruction is the PPN-\(\gamma\) obstruction]
\label{prop:trace-obstruction-gamma}
In the isotropic two-sector weak-field regime, suppose the logarithmic cell dilation has the expansion
\[
        u(r)=\zeta_C\frac{\Phi_C(r)}{c^2}+O(c^{-4}).
\]
Then
\[
        \gamma_C=1-\zeta_C.
\]
Thus the GR spatial coefficient \(\gamma_C=1\) is equivalent to vanishing first-order trace dilation \(\zeta_C=0\).
\end{proposition}

\begin{proof}
By definition,
\[
        e^{2u}=\Theta_t\Theta_s.
\]
Using the two-sector weak-field expansion,
\[
\begin{aligned}
        \Theta_t\Theta_s
        &=\left(1+2\frac{\Phi_C}{c^2}\right)
          \left(1-2\gamma_C\frac{\Phi_C}{c^2}\right)+O(c^{-4})\\
        &=1+2(1-\gamma_C)\frac{\Phi_C}{c^2}+O(c^{-4}).
\end{aligned}
\]
Meanwhile \(e^{2u}=1+2\zeta_C\Phi_C/c^2+O(c^{-4})\).  Equating coefficients gives \(\zeta_C=1-\gamma_C\), hence \(\gamma_C=1-\zeta_C\).
\end{proof}

\begin{audit}[The remaining obstruction]
The spatial gate has now been reduced to a concrete question: does static scalar geometric recruitment source the trace mode \(u_e\) of the local clock-ruler response?  If not, fixed-cell variational response gives trace-free transport and closes \(\gamma_C=1\).  If yes, the sourced trace coefficient \(\zeta_C\) shifts the PPN spatial coefficient by \(\gamma_C=1-\zeta_C\).  The obstruction is no longer vague.  It is the existence of a first-order \(\chi u_e\) term in the effective local response potential.
\end{audit}


\subsection{Generator form: the obstruction is a \texorpdfstring{$GL(2)$}{GL(2)} trace}

The variational statement above is useful, but the same obstruction can be stated more invariantly.  A weak static load acts locally on each clock-ruler comparison plane by a linear response map.  The determinant of this map is the comparison-cell factor.  The trace of its infinitesimal generator is therefore exactly the first-order cell-dilation obstruction.

\begin{definition}[Clock-ruler response map]
\label{def:clock-ruler-response-map}
Fix a represented spatial direction \(e\).  After removing coordinate relabelings and shear, write the loaded clock-ruler increments as
\[
        \begin{pmatrix} d\tau \\ d\ell_e \end{pmatrix}
        =R_e(\hat\chi)
        \begin{pmatrix} d\tau_0 \\ d\ell_{e0} \end{pmatrix},
        \qquad
        \hat\chi(r):=\frac{\Phi_C(r)}{c^2}=O(c^{-2}),
\]
where \(R_e(0)=I\) and \(R_e(\hat\chi)\in GL^+(2,\mathbb R)\) in the represented smooth regime.  Its infinitesimal generator is
\[
        L_e:=\left.\frac{dR_e}{d\hat\chi}\right|_{0},
        \qquad
        R_e(\hat\chi)=I+\hat\chi L_e+O(\hat\chi^2).
\]
The trace component
\[
        \theta_e:=\operatorname{tr}L_e
\]
is the infinitesimal comparison-cell dilation rate.  The traceless part \(L_e-\frac12\theta_e I\) is the exchange/shear part of the clock-ruler response.
\end{definition}

\begin{proposition}[Generator trace equals trace obstruction]
\label{prop:generator-trace-obstruction}
In the diagonal isotropic weak-field sector, the logarithmic cell dilation satisfies
\[
        u_e(r)=\theta_e\frac{\Phi_C(r)}{c^2}+O(c^{-4}).
\]
Consequently
\[
        \zeta_C=\theta_e,
        \qquad
        \gamma_C=1-\theta_e.
\]
Thus the GR spatial coefficient \(\gamma_C=1\) is equivalent, at this order, to the infinitesimal load generator being traceless:
\[
        \operatorname{tr}L_e=0.
\]
\end{proposition}

\begin{proof}
The comparison-cell factor is
\[
        \mathcal A_e=\det R_e(\hat\chi)
        =1+\hat\chi\operatorname{tr}L_e+O(\hat\chi^2).
\]
Since \(u_e=\log\mathcal A_e\), we obtain
\[
        u_e=\hat\chi\operatorname{tr}L_e+O(\hat\chi^2)
        =\theta_e\Phi_C/c^2+O(c^{-4}).
\]
Comparison with \Cref{prop:trace-obstruction-gamma} gives \(\zeta_C=\theta_e\) and hence \(\gamma_C=1-\theta_e\).
\end{proof}


\subsection{Gauge-reduced scalar normal form}

The generator formulation still permits too many apparent degrees of freedom.  Some are physical, but some are only coordinate choice, anisotropic shear, or time-space mixing inappropriate to the static scalar exterior sector.  The next reduction is to put the local response into normal form.

\begin{definition}[Gauge-reduced scalar response sector]
\label{def:gauge-reduced-scalar-sector}
Fix a static, spherically symmetric exterior weak-field load.  A local clock-ruler response map is in the gauge-reduced scalar sector when, after quotienting coordinate relabelings, anisotropic shear, and parity-forbidden time-space mixing, its infinitesimal generator is diagonal and direction-independent:
\[
        L_{\rm sc}
        =
        \begin{pmatrix}
        \lambda_t & 0\\
        0 & \lambda_s
        \end{pmatrix}
        =\frac{\theta}{2}I+\sigma H,
        \qquad
        H:=
        \begin{pmatrix}
        1&0\\
        0&-1
        \end{pmatrix}.
\]
Here
\[
        \theta:=\lambda_t+\lambda_s=\operatorname{tr}L_{\rm sc},
        \qquad
        \sigma:=\frac{\lambda_t-\lambda_s}{2}.
\]
The trace \(\theta\) is the comparison-cell dilation mode; the exchange coefficient \(\sigma\) transfers comparison scale between clock and ruler channels without changing the cell determinant.
\end{definition}

\begin{proposition}[Clock calibration leaves only the trace obstruction]
\label{prop:normal-form-trace-obstruction}
In the gauge-reduced scalar sector, suppose the clock response is calibrated by the redshift recovery,
\[
        \sqrt{\Theta_t}=1+\frac{\Phi_C}{c^2}+O(c^{-4}),
\]
and write the spatial response as
\[
        \sqrt{\Theta_s}=1-\gamma_C\frac{\Phi_C}{c^2}+O(c^{-4}).
\]
Then the normal-form generator has
\[
        \lambda_t=1,
        \qquad
        \lambda_s=-\gamma_C,
\]
so that
\[
        \theta=1-
        \gamma_C,
        \qquad
        \sigma=\frac{1+
        \gamma_C}{2}.
\]
Therefore, once the clock sector is calibrated, the entire first-order spatial-response problem is the trace problem:
\[
        \gamma_C=1
        \quad\Longleftrightarrow\quad
        \theta=0.
\]
\end{proposition}

\begin{proof}
By \Cref{def:clock-ruler-response-map}, the loaded increments are
\[
        d\tau=(1+\hat\chi\lambda_t)d\tau_0+O(\hat\chi^2),
        \qquad
        d\ell=(1+\hat\chi\lambda_s)d\ell_0+O(\hat\chi^2),
\]
with \(\hat\chi=\Phi_C/c^2\).  Matching these expressions to the displayed clock and ruler factors gives \(\lambda_t=1\) and \(\lambda_s=-\gamma_C\).  The formulae for \(\theta\) and \(\sigma\) follow from \Cref{def:gauge-reduced-scalar-sector}.  Hence \(\gamma_C=1\) iff \(\theta=0\).
\end{proof}

\begin{corollary}[Same-sign depletion failure in normal form]
\label{cor:same-sign-failure-normal-form}
A same-sign scalar depletion of clock and ruler increments has \(\lambda_t=\lambda_s=1\).  With the clock sector calibrated, this implies
\[
        \gamma_C=-1,
\]
so it reproduces the redshift sign but fails the spatial weak-field sector.
\end{corollary}

\begin{proof}
If \(\lambda_t=1\) and \(\lambda_s=1\), then \(-\gamma_C=1\), hence \(\gamma_C=-1\).  Equivalently, the trace is \(\theta=2\), so \(\gamma_C=1-\theta=-1\).
\end{proof}

\begin{proposition}[Unique scalar closure compatible with calibrated weak-field GR]
\label{prop:unique-scalar-closure}
Within the gauge-reduced scalar sector with the clock coefficient fixed by redshift, the following are equivalent to first post-Newtonian order:
\[
\begin{aligned}
        L_{\rm sc}\in\mathfrak{sl}(2,\mathbb R)
        &\Longleftrightarrow
        L_{\rm sc}=H \\
        &\Longleftrightarrow
        R_{\rm sc}\in SL(2,\mathbb R)+O(c^{-4}) \\
        &\Longleftrightarrow
        \Theta_s=\Theta_t^{-1}+O(c^{-4}) \\
        &\Longleftrightarrow
        \gamma_C=1.
\end{aligned}
\]
Thus the scalar weak-field closure is unique: the load must act as a trace-free clock-ruler exchange, not as a common dilation.
\end{proposition}

\begin{proof}
Clock calibration gives \(\lambda_t=1\).  Tracelessness gives \(\lambda_s=-1\), hence \(L_{\rm sc}=H\).  This is equivalent to \(\det R_{\rm sc}=1+O(c^{-4})\), since the determinant changes to first order by \(\operatorname{tr}L_{\rm sc}\hat\chi\).  The determinant condition is \(\Theta_s\Theta_t=1+O(c^{-4})\), which is the same as \(\Theta_s=\Theta_t^{-1}+O(c^{-4})\).  By \Cref{prop:normal-form-trace-obstruction}, this is equivalent to \(\gamma_C=1\).
\end{proof}

\begin{theorem}[Scalar weak-field closure package]
\label{thm:scalar-closure-package}
Assume the static isotropic exterior weak-field sector, the metric-representation gate, universal matter coupling to the same comparison kernel, and the clock calibration supplied by the redshift recovery.  After coordinate gauge, anisotropic shear, and parity-forbidden time-space mixing are removed, the following statements are equivalent to first post-Newtonian order:
\[
\begin{aligned}
        \text{no scalar comparison-cell dilation}
        &\Longleftrightarrow
        \operatorname{tr}L_{\rm sc}=0 \\
        &\Longleftrightarrow
        L_{\rm sc}\in\mathfrak{sl}(2,\mathbb R) \\
        &\Longleftrightarrow
        \det R_{\rm sc}=1+O(c^{-4}) \\
        &\Longleftrightarrow
        \Theta_s\Theta_t=1+O(c^{-4}) \\
        &\Longleftrightarrow
        \gamma_C=1.
\end{aligned}
\]
If these equivalent conditions hold, the conditional weak-field exterior recovery theorem applies.  If they fail with trace \(\theta\), then
\[
        \gamma_C=1-\theta,
\]
so two-sector weak-field observations measure the scalar comparison-cell dilation directly.
\end{theorem}

\begin{proof}
Clock calibration gives \(\lambda_t=1\).  In the gauge-reduced scalar sector, \Cref{prop:normal-form-trace-obstruction} gives \(\theta=1-\gamma_C\), so \(\gamma_C=1\) iff \(\theta=0\).  The equivalence between vanishing trace, membership in \(\mathfrak{sl}(2,\mathbb R)\), first-order unimodularity of \(R_{\rm sc}\), and preservation of the normalized clock-ruler cell is the determinant identity for a first-order \(GL(2)\) response.  In the isotropic two-sector kernel this determinant condition is \(\Theta_s\Theta_t=1+O(c^{-4})\).  The conditional weak-field recovery then follows from \Cref{thm:conditional-weak-field}.  If \(\theta\neq0\), \Cref{prop:trace-obstruction-gamma} gives the shifted coefficient \(\gamma_C=1-\theta\).
\end{proof}

\begin{audit}[What the scalar package reduces]
The spatial problem is no longer an unconstrained metric-fitting problem.  In the static scalar exterior sector, after gauge and shear are removed, redshift fixes \(\lambda_t\).  The only remaining first-order scalar ambiguity is \(\theta=\operatorname{tr}L_{\rm sc}\).  Proving the APF substrate has no independent scalar comparison-cell dilation is therefore enough to close the leading two-sector weak-field geometry.
\end{audit}

\begin{gate}[Unimodular scalar load response]
\label{gate:unimodular-load-response}
Static scalar geometric recruitment acts on each local clock-ruler comparison plane by a unimodular response after coordinate gauge, anisotropic shear, and composition-dependent channels are removed:
\[
        \det R_e(\hat\chi)=1+O(c^{-4}),
        \qquad\text{equivalently}\qquad
        L_e\in\mathfrak{sl}(2,\mathbb R).
\]
Physical reading: scalar recruitment may redistribute comparison capacity between clock and ruler channels, but it does not create or destroy the elementary local comparison cell.  A failure of this gate is not a small technical issue; it is an additional scalar trace mode in the weak-field geometry.
\end{gate}

\begin{proposition}[Unimodular response closes the spatial coefficient]
\label{prop:unimodular-implies-gamma}
Assume the two-sector weak-field kernel and the unimodular scalar load response gate.  Then
\[
        \Theta_s\Theta_t=1+O(c^{-4}),
        \qquad
        \gamma_C=1.
\]
\end{proposition}

\begin{proof}
The determinant of the clock-ruler response map is the normalized comparison-cell factor:
\[
        \det R_e=\mathcal A_e
        =\left(\frac{d\tau}{dt}\right)
         \left(\frac{d\ell_e}{dx_e}\right)
        =\sqrt{\Theta_t\Theta_{s,e}}.
\]
Unimodularity gives \(\mathcal A_e=1+O(c^{-4})\), hence \(\Theta_t\Theta_{s,e}=1+O(c^{-4})\).  In the isotropic sector this is \(\Theta_t\Theta_s=1+O(c^{-4})\), and \Cref{prop:reciprocal-implies-gamma} gives \(\gamma_C=1\).
\end{proof}

\begin{corollary}[Equivalence of the spatial closure targets]
\label{cor:spatial-target-equivalence}
In the isotropic weak-field regime, the following are equivalent to first post-Newtonian order:
\[
\begin{aligned}
        \operatorname{tr}L_e=0
        &\Longleftrightarrow
        \det R_e=1+O(c^{-4}) \\
        &\Longleftrightarrow
        \mathcal A_e=1+O(c^{-4}) \\
        &\Longleftrightarrow
        \Theta_s\Theta_t=1+O(c^{-4}) \\
        &\Longleftrightarrow
        \gamma_C=1.
\end{aligned}
\]
The spatial gate can therefore be pursued as any one of four equivalent problems: rule out the trace generator, prove unimodular response, prove comparison-cell conservation, or derive reciprocal clock-ruler response.
\end{corollary}

\begin{audit}[Operational diagnostic for the trace mode]
If weak-field tests determine an effective PPN coefficient \(\gamma_{\rm obs}\), the APF trace mode is read off as
\[
        \theta_e=\zeta_C=1-\gamma_{\rm obs}.
\]
Thus the observational role of light bending, Shapiro delay, and other two-sector probes is not merely to test GR.  In the cost language, they bound or detect the first-order determinant of the local clock-ruler response map.  Redshift alone cannot do this, because it probes only the clock projection.
\end{audit}

\begin{proposition}[Nonzero trace is a scalar comparison mode]
\label{prop:nonzero-trace-scalar-mode}
Within the two-sector weak-field ansatz, a nonzero first-order trace generator \(\theta_e\neq0\) is observationally equivalent, at this order, to adding an independent scalar comparison mode that changes the spatial PPN coefficient while leaving the redshift sector fixed.
\end{proposition}

\begin{proof}
The redshift sector fixes \(\Theta_t=1+2\Phi_C/c^2+O(c^{-4})\).  A nonzero trace generator changes the product \(\Theta_t\Theta_s\) by
\[
        \Theta_t\Theta_s
        =1+2\theta_e\frac{\Phi_C}{c^2}+O(c^{-4}),
\]
which by \Cref{prop:generator-trace-obstruction} shifts \(\gamma_C\) to \(1-\theta_e\).  Since \(\Theta_t\) is unchanged, the new degree of freedom is invisible to clock-sector redshift but visible to two-sector spatial observables.  That is exactly the weak-field role of an additional scalar comparison mode.
\end{proof}

\begin{remark}[Why this is a strengthening]
The variational obstruction can be expressed invariantly as a generator question: is scalar load transport \(GL(2)\) with nonzero trace, or \(SL(2)\) to first order?  This is the cleanest current statement of the spatial-response problem.
\end{remark}

\begin{corollary}[A sharper spatial-response target]
\label{cor:sharper-target}
The weak-field spatial sector can be closed by proving the unimodular scalar load response gate \Cref{gate:unimodular-load-response}, or equivalently any of the targets listed in \Cref{cor:spatial-target-equivalence}.  One convenient proof chain is
\[
\begin{aligned}
        \operatorname{tr}L_e=0
        &\Longrightarrow
        \det R_e=1+O(c^{-4}) \\
        &\Longrightarrow
        \Theta_s\Theta_t=1+O(c^{-4}) \\
        &\Longrightarrow
        \gamma_C=1 \\
        &\Longrightarrow
        \text{leading two-sector weak-field GR recovery}.
\end{aligned}
\]
Thus the next target is no longer the bare coefficient \(\gamma_C\), but the underlying no-trace / unimodular comparison-load certificate.
\end{corollary}

\begin{assumption}[Two-sector weak-field response]
\label{ass:two-sector-response}
In the exterior weak-field regime generated by a static point source, the represented comparison kernel has
\[
        \Theta_t(r)=1+\frac{2\Phi_C(r)}{c^2}+O(c^{-4}),
        \qquad
        \Theta_s(r)=1-2\gamma_C\frac{\Phi_C(r)}{c^2}+O(c^{-4}).
\]
The clock calculation gives the first relation conditionally.  The second relation is the spatial-response target.
\end{assumption}

\begin{theorem}[Conditional weak-field exterior recovery]
\label{thm:conditional-weak-field}
Assume the metric-representation gate, the exterior point-source profile, the Newtonian calibration \(\Phi_C=-GM/r\), the proper-time dictionary, universal matter coupling to the same comparison kernel, and the spatial-response gate \(\gamma_C=1\).  Then the represented exterior weak-field cost kernel is
\[
        ds_C^2
        =-\left(1-\frac{2GM}{rc^2}\right)c^2dt^2
        +\left(1+\frac{2GM}{rc^2}\right)d\mathbf x^2
        +O(c^{-4}),
\]
which is the first-post-Newtonian Schwarzschild exterior metric in isotropic weak-field form.  Consequently the framework conditionally recovers, at leading weak-field order, gravitational redshift, clock-height tests, Newtonian acceleration, light bending, and Shapiro time delay.
\end{theorem}

\begin{proof}
The clock-sector recovery and calibration give
\[
        \Theta_t(r)=1-\frac{2GM}{rc^2}+O(c^{-4}).
\]
The spatial-response gate with \(\gamma_C=1\) gives
\[
        \Theta_s(r)=1+\frac{2GM}{rc^2}+O(c^{-4}).
\]
Substitution into \Cref{def:two-sector-kernel} gives the displayed metric kernel.  In standard post-Newtonian bookkeeping, this pair of coefficients is exactly the weak-field Schwarzschild exterior through the leading order controlling the listed observables.  The recovery is conditional because the metric-representation gate, universal matter coupling, and either the spatial-response gate or its sharper reciprocal clock-ruler route are not yet derived from substrate primitives.
\end{proof}


\begin{corollary}[No-trace scalar response is sufficient for the leading weak-field package]
\label{cor:no-trace-suffices-weak-field}
Under the hypotheses of \Cref{thm:conditional-weak-field}, the explicit spatial-response assumption \(\gamma_C=1\) may be replaced by any one of the equivalent scalar-closure conditions in \Cref{thm:scalar-closure-package}.  In particular, if the gauge-reduced scalar load generator is traceless,
\[
        \operatorname{tr}L_{\rm sc}=0,
\]
then the conditional first-post-Newtonian exterior recovery follows.
\end{corollary}

\begin{proof}
By \Cref{thm:scalar-closure-package}, vanishing trace is equivalent to \(\gamma_C=1\) in the static isotropic scalar exterior sector after clock calibration.  Substituting this equivalent gate for the spatial-response assumption gives exactly the hypotheses of \Cref{thm:conditional-weak-field}.
\end{proof}

\begin{corollary}[Observable dependency split]
\label{cor:observable-dependency-split}
Within the two-sector weak-field kernel, redshift and clock-height tests are clock-sector observables, while light bending and Shapiro delay are two-sector observables.  Therefore any APF claim to recover leading weak-field GR must close both the clock-response gate and the spatial-response gate.
\end{corollary}

\begin{audit}[No-smuggling table for the weak-field suite]
\begin{longtable}{p{0.22\textwidth}p{0.16\textwidth}p{0.16\textwidth}p{0.33\textwidth}}
\toprule
Observable & Needs \(\Theta_t\) & Needs \(\Theta_s\) & Present status \\
\midrule
Static redshift & Yes & No & Recovered conditionally from clock response. \\
Clock-height tests & Yes & No & Same gate as redshift, with local calibration. \\
Newtonian acceleration & Yes & No at leading slow-body order & Conditional on geodesic-cost dictionary and universal matter coupling. \\
Light bending & Yes & Yes & Open until \(\gamma_C=1\) is derived. \\
Shapiro delay & Yes & Yes & Open until null comparison uses the two-sector kernel with \(\gamma_C=1\). \\
Perihelion precession & Yes & Yes plus next-order terms & Partly first-PN; full recovery needs nonlinear corrections. \\
Frame dragging & Needs \(h_{0i}\) & Yes & Deferred to rotating-source tensor sector. \\
Gravitational radiation & Tensor dynamics & Tensor dynamics & Deferred to dynamical self-coupled geometric substrate. \\
\bottomrule
\end{longtable}
\end{audit}

\begin{remark}[Strategic consequence]
The next decisive mathematical target is the spatial-response lemma: derive \(\gamma_C=1\) from the same substrate principles that produce the clock-stiffness depletion.  Closing that lemma would upgrade this draft from a redshift proof-of-principle to a genuine weak-field GR recovery program.  The scalar package sharpens the target again: prove that the scalar load response generator is traceless, equivalently \(R_e\in SL(2,\mathbb R)\) to first order, or compute the trace obstruction.  By \Cref{prop:unimodular-implies-gamma,cor:spatial-target-equivalence}, that certificate closes the leading spatial coefficient.  By \Cref{prop:generator-trace-obstruction,prop:trace-obstruction-gamma}, any trace source is exactly the PPN-\(\gamma\) obstruction.
\end{remark}

\section{From scalar clock response to tensorial comparison response}

The scalar calculation used one effective load \(\chi\) and one stiffness \(K_t\).  A geometric theory requires a tensorial comparison kernel.  Let \(Q_{ab}(x)\) denote the local quadratic cost kernel for infinitesimal comparison continuations.  In a represented smooth regime, \(Q_{ab}\) is read as a metric kernel up to normalization and signature convention.

\begin{definition}[Tensorial comparison kernel]
A tensorial comparison kernel is a smooth field \(Q_{ab}\) such that the leading squared cost of an infinitesimal comparison displacement \(dx^a\) is
\[
        d\ell_C^2 = Q_{ab}(x)\,dx^a dx^b.
\]
When the signature gate closes, \(Q_{ab}\) has Lorentzian signature and is represented as a spacetime metric kernel.
\end{definition}

\begin{definition}[Weak perturbation around unloaded substrate]
In the weak-field regime, write
\[
        Q_{ab}=\eta_{ab}+h_{ab},
        \qquad |h_{ab}|\ll 1,
\]
where \(\eta_{ab}\) is the unloaded Lorentzian comparison kernel and \(h_{ab}\) is the recruited comparison-cost distortion.
\end{definition}

\begin{principle}[Universal comparison-load coupling]
Matter distinctions couple to geometry only through the comparison burden they impose on the substrate.  In the represented tensorial regime this burden is summarized by a stress-energy functional \(T^{ab}\), and the leading coupling has the schematic form
\[
        \Erec[h;T]
        =\Egeom[h]-\lambda\int h_{ab}T^{ab}\,d^4x+O(h^2T,hT^2).
\]
Universality of free fall requires the same \(h_{ab}\) to couple to all admissible matter sectors through the same comparison-load channel.
\end{principle}

\begin{gate}[No composition dependence gate]
The equivalence-principle-like claim closes only if the recruitment load of a composite body depends on its total carried stress-energy distinction and not on its internal microscopic composition, except through already-accounted internal energy.  In APF terms, all admissible clocks and test bodies must read the same geometric comparison kernel.
\end{gate}

\begin{proposition}[Redshift sector as the \(h_{00}\) projection]
In the weak static regime, the redshift recovery fixes the projection
\[
        h_{00}(r)\simeq \frac{2\Phi_C(r)}{c^2}=-\frac{2GM}{rc^2}
\]
up to sign convention for the metric kernel.  It does not fix \(h_{ij}\) or \(h_{0i}\).
\end{proposition}

\begin{proof}
The clock factor is determined by the time-time component of the quadratic comparison kernel.  Static redshift compares only rest-frame time increments, so it projects onto the time-time component.  Spatial and mixed components require spatial and moving-probe comparisons.
\end{proof}

\section{Translation route to the Paper 6 field equations}

The Paper~6 route to GR is not to guess the Schwarzschild metric, but to characterize the admissible local cost functional for \(Q_{ab}\) after the represented smooth Lorentzian gates close.  The expected structure is:
\[
        \Egeom[Q]
        =\int_{\Mgeom} \Bigl[\text{local scalar made from }Q\text{ and its first/second derivatives}\Bigr]\,d\mu_Q.
\]
Paper~6 supplies this represented-geometry closure using H1--H3, A9, and Lovelock uniqueness.  The Paper~9 challenge is narrower: translate those continuum hypotheses into finite-continuability certificates and identify their weak-field scalar face.

\begin{gate}[Coordinate-relabeling invariance gate]
The cost assigned to a comparison-continuation configuration may not depend on arbitrary labels used for comparison sites.  In the smooth represented regime, this becomes diffeomorphism-type invariance of the geometric cost functional.
\end{gate}

\begin{gate}[Second-order locality gate]
If the substrate cost is local and the represented field equations are to avoid additional higher-derivative propagating comparison modes, the continuum Euler-Lagrange equations should be second order in the metric kernel.  This is the gate that points toward curvature scalar functionals rather than arbitrary higher-curvature actions.
\end{gate}

\begin{gate}[Bianchi/ledger gate]
The geometric side of the represented field equation must have an identically conserved ledger divergence so that local stress-energy continuation is compatible with the geometry equation.  In GR this is the contracted Bianchi identity.  In cost language it is the statement that relabeling comparison sites cannot create or destroy committed load.
\end{gate}


\begin{definition}[Geometric Euler ledger]
Let \(\Egeom[Q]\) be a represented geometric cost functional.  Its Euler ledger is the variational tensor density \(\mathcal E^{ab}\) defined by
\[
        \delta \Egeom[Q]
        =\int \mathcal E^{ab}\,\delta Q_{ab}\,d\mu_Q
        +\text{boundary terms}.
\]
A matter recruitment functional contributes a ledger \(\mathcal T^{ab}\) in the same variational channel.
\end{definition}

\begin{proposition}[Relabeling invariance implies a conserved geometric ledger]
\label{prop:ledger-identity}
If \(\Egeom[Q]\) is invariant under compactly supported infinitesimal relabelings of comparison sites, then its Euler ledger satisfies a formal conservation identity
\[
        \nabla_a\mathcal E^{ab}=0
\]
in the represented smooth metric regime, up to the boundary conventions of the functional.
\end{proposition}

\begin{proof}[Proof sketch]
An infinitesimal relabeling generated by a vector field \(\xi^a\) changes the metric kernel by its Lie derivative,
\[
        \delta Q_{ab}=\Lie_\xi Q_{ab}=\nabla_a\xi_b+\nabla_b\xi_a.
\]
Relabeling invariance gives \(\delta\Egeom=0\) for all compactly supported \(\xi\).  Substituting the variation and integrating by parts yields
\[
        0=-2\int (\nabla_a\mathcal E^{ab})\xi_b\,d\mu_Q.
\]
Since \(\xi\) is arbitrary, \(\nabla_a\mathcal E^{ab}=0\).  The result is a represented-regime conservation identity; translating it back to finite continuations is a separate APF gate.
\end{proof}

\begin{corollary}[Matter compatibility condition]
If the represented field equation has the form
\[
        \mathcal E^{ab}=\lambda \mathcal T^{ab},
\]
then relabeling invariance of the geometric ledger requires
\[
        \nabla_a\mathcal T^{ab}=0
\]
inside the same represented regime.  In APF terms, local committed recruitment load must continue consistently with the comparison geometry that carries it.
\end{corollary}

\begin{theorem}[Translation target: Einstein-Hilbert uniqueness, schematic]
\label{thm:target-eh}
If the represented geometric cost functional is local, smooth, coordinate-relabeling invariant, built only from the tensorial comparison kernel and its derivatives, yields second-order equations, and has a conserved ledger divergence compatible with universal stress-energy recruitment, then the leading continuum functional is of Einstein-Hilbert type,
\[
        \Egeom[Q]\propto \int (R[Q]-2\Lambda)\,d\mu_Q,
\]
up to boundary terms, normalization, and higher-order corrections outside the gate.
\end{theorem}

\begin{warning}[Status of \Cref{thm:target-eh}]
This theorem is not reproved here.  It is the Paper~9 translation target corresponding to the Paper~6 A9/Lovelock closure.  Paper~6 supplies the represented-geometry proof; Paper~9 asks which finite-continuability certificates force the same continuum hypotheses and, in the weak-field scalar sector, which certificate forces the no-trace condition \(\theta=0\).
\end{warning}

\section{Nonlinear alignment: from response law to the Paper 6 closure}

The leading-order theory uses a scalar accumulated load \(\chi\).  Paper~6 supplies the represented tensorial, nonlinear, self-coupled closure to the Einstein equations.  The remaining Paper~9 problem is to show how the finite comparison-substrate response lifts into exactly the Paper~6 tensorial closure rather than an alternative scalar-tensor or higher-derivative geometry.

\begin{gate}[Tensorial geometric response gate]
The scalar stiffness \(K_t\) must lift to a tensorial comparison-cost structure \(g_{ab}\) or equivalent cost kernel such that matter recruitment load changes the full comparison geometry, not only the static time-comparison sector.
\end{gate}

\begin{gate}[Self-coupling gate]
The geometric recruitment load must include the cost of the geometric substrate's own recruited configuration.  In GR language, gravitational field energy gravitates.  In cost language, committed comparison-cost distortion contributes to the future comparison-cost kernel.
\end{gate}

\begin{gate}[Einstein-equation gate]
The nonlinear geometric recruitment functional must have Euler-Lagrange equations equivalent, in the represented smooth Lorentzian regime, to
\[
        G_{ab}+\Lambda g_{ab}=\frac{8\pi G}{c^4}T_{ab},
\]
with \(G\) remaining either a substrate-specific input or a derived stiffness constant.
\end{gate}

\begin{warning}[Primary translation theorem]
The Paper~6 endpoint can be restated in Paper~9 language as follows:

\emph{If finite comparison continuations admit a smooth local Lorentzian quadratic kernel, universal recruitment coupling, local A9-type closure, and nonlinear self-consistent cost response, then their represented continuum equations are the Einstein equations.}

Paper~6 proves the represented-geometry closure under its hypotheses.  Paper~9 has not yet proved that finite continuability forces those hypotheses.  Its current hard local target is the weak-field trace theorem \(\theta=0\), the scalar shadow of the same closure.
\end{warning}

\section{Input ledger and closure plan}

\begin{longtable}{p{0.27\textwidth}p{0.12\textwidth}p{0.52\textwidth}}
\toprule
Claim / object & Status & What remains \\
\midrule
Continuation judgment \(\G\entails_C D\goes E\) & \Status{P} & Imported from Paper 10 as primitive grammar. \\
Comparison continuation & \Status{P/C} & Defined here; experimental/operational certification still needs examples. \\
Metric representation from comparison cost & \Status{C} & Conditional on local completeness, smoothness, quadratic tangent limit, and path additivity. \\
Lorentzian signature & \Status{O} & Requires finite-coordination cone plus irreversible record orientation proof. \\
Four-dimensional lock-in & \Status{O} & Requires capacity/mode-count argument unified with signature. \\
Quadratic recruitment functional & \Status{C} & Assumed as linear-response effective form. \\
Point-source \(1/r\) profile & \Status{C/P} & Proved from assumed isotropic quadratic kernel; kernel itself not derived. \\
Cost-curvature response \(K_t=K_{t0}(1-\alpha\chi)\) & \Status{C} & Needs perturbative derivation from nonlinear recruitment functional. \\
Proper-time dictionary & \Status{C/R} & Representation dictionary; needs derivation from identity-continuation action. \\
Redshift recovery & \Status{R} & Recovered under response law, profile, dictionary, and \(G\)-calibration. \\
Two-sector weak-field kernel & \Status{C} & Defined here as the weak-field split between clock and spatial comparison sectors; the closure route proceeds through reciprocal response, cell conservation, trace-free transport, fixed-cell variational response, and finally the gauge-reduced trace normal form. \\
Fixed-cell variational response & \Status{O/C} & Sharper candidate certificate; if proved, implies trace-free transport/cell conservation by \Cref{prop:fixed-cell-implies-trace-free}. \\
Canonical trace-free clock-ruler transport & \Status{O/C} & Candidate certificate; follows from fixed-cell variational response and implies cell conservation by \Cref{prop:trace-free-implies-cell}. \\
Trace obstruction coefficient & \Status{P/C} & By \Cref{prop:generator-trace-obstruction,prop:trace-obstruction-gamma}, a first-order trace generator is exactly the obstruction to \(\gamma_C=1\). \\
Gauge-reduced scalar normal form & \Status{P/C} & After clock calibration, the only first-order scalar ambiguity is the generator trace by \Cref{prop:normal-form-trace-obstruction,thm:scalar-closure-package}. \\
Unimodular scalar load response & \Status{O/C} & Closure target sharpened by the normal form; if proved, closes the spatial coefficient by \Cref{prop:unimodular-implies-gamma,prop:unique-scalar-closure,thm:scalar-closure-package}. \\
Directionwise clock-ruler cell conservation & \Status{O/C} & Candidate operational certificate; follows from trace-free transport by \Cref{prop:trace-free-implies-cell}; implies reciprocal response by \Cref{prop:cell-implies-reciprocal}. \\
Reciprocal clock-ruler response \(\Theta_s\Theta_t=1\) & \Status{C/P} & Follows from cell conservation in the isotropic two-sector regime; implies \(\gamma_C=1\). \\
Spatial response \(\gamma_C=1\) & \Status{C/O} & Follows from reciprocal response by \Cref{prop:reciprocal-implies-gamma}; reciprocal gate itself remains open. \\
Conditional weak-field exterior recovery & \Status{C/R} & Follows if clock response, spatial response, and universal coupling close. \\
Newtonian acceleration & \Status{C/R} & Recovered under geodesic-cost dictionary. \\
Newton's constant \(G\) & \Status{O/Input} & Treated as substrate stiffness input; deriving value is open. \\
Full Einstein equations & \Status{I/C} & Imported from Paper~6 under H1--H3/A9/Lovelock; Paper~9 translation back to finite-continuability certificates remains open. \\
\bottomrule
\end{longtable}

\section{Formal theorem tree for closure}

The following theorem tree is the working spine for later versions.  It separates representation theorems from physical recovery theorems.

\begin{longtable}{p{0.09\textwidth}p{0.30\textwidth}p{0.10\textwidth}p{0.36\textwidth}}
\toprule
ID & Claim & Status & Dependency \\
\midrule
T26.1 & Finite comparison continuations define a local cost germ. & \Status{C} & Existence of refinement limits. \\
T26.2 & Regular cost germs define Finsler-type comparison length. & \Status{C/P} & Convexity and path-additivity gates. \\
T26.3 & Parallelogram certificate gives a quadratic metric kernel. & \Status{P} & \Cref{prop:parallelogram-kernel}. \\
T26.4 & Finite coordination plus irreversible records give Lorentzian signature. & \Status{O} & Signature gate. \\
T26.5 & Capacity/mode-counting locks the represented dimension to \((1+3)\). & \Status{O} & Dimension gate. \\
T26.6 & Isotropic quadratic exterior recruitment gives \(\chi\sim1/r\). & \Status{C/P} & Assumed exterior kernel; \Cref{prop:point-source-profile}. \\
T26.7 & Bilinear load-probe coupling gives linear clock-stiffness depletion. & \Status{C/P} & Coupling assumption; \Cref{prop:response-from-coupling}. \\
T26.8 & Calibrated clock-stiffness depletion recovers gravitational redshift. & \Status{R} & Proper-time dictionary and \(G\)-calibration. \\
T26.8a & Fixed-cell variational response implies cell conservation. & \Status{C/P} & \Cref{prop:fixed-cell-implies-trace-free}; proving the variational gate is open. \\
T26.8b & Trace-free clock-ruler load transport implies cell conservation. & \Status{C/P} & \Cref{prop:trace-free-implies-cell}; follows if fixed-cell response closes. \\
T26.8c & Trace obstruction is the PPN-\(\gamma\) obstruction. & \Status{C/P} & \Cref{prop:trace-obstruction-gamma}. \\
T26.8d & Generator trace equals the trace obstruction. & \Status{C/P} & \Cref{prop:generator-trace-obstruction}. \\
T26.8e & Unimodular scalar load response implies \(\gamma_C=1\). & \Status{C/P} & \Cref{prop:unimodular-implies-gamma}; proving the gate is open. \\
T26.8f & Gauge-reduced scalar normal form makes the trace the only remaining first-order scalar obstruction. & \Status{P/C} & \Cref{prop:normal-form-trace-obstruction,prop:unique-scalar-closure}. \\
T26.8g & Scalar closure package reduces weak-field spatial recovery to no-trace/unimodular load response. & \Status{P/C} & \Cref{thm:scalar-closure-package}. \\
T26.9 & Directionwise clock-ruler cell conservation implies reciprocal response. & \Status{C/P} & \Cref{prop:cell-implies-reciprocal}; proving the cell-conservation gate is open. \\
T26.9a & Reciprocal clock-ruler response implies \(\gamma_C=1\). & \Status{C/P} & \Cref{prop:reciprocal-implies-gamma}. \\
T26.9b & Spatial response fixes \(\gamma_C=\gamma_{\rm PPN}=1\). & \Status{C/O} & Follows if \Cref{gate:cell-conservation} closes. \\
T26.10 & Conditional two-sector weak-field recovery gives the first-PN Schwarzschild exterior. & \Status{C/R} & Clock response plus T26.9; \Cref{thm:conditional-weak-field}. \\
T26.11 & Relabeling invariance gives a conserved geometric ledger. & \Status{C/P} & Represented smooth regime; \Cref{prop:ledger-identity}. \\
T26.12 & Local second-order self-coupled ledger gives Einstein equations. & \Status{I/C} & Paper~6 A9/Lovelock import; finite-continuability translation remains open. \\
\bottomrule
\end{longtable}


\section{Minimal publishable core}
\label{sec:minimal-publishable-core}

A stable version of this paper should not try to close the full nonlinear theory.  The minimal publishable theorem package is smaller:
\begin{enumerate}
\item a clean comparison-cost-to-metric representation lemma;
\item a conditional redshift recovery from calibrated clock response;
\item a proof that clock observables alone do not determine spatial response;
\item the scalar normal-form theorem reducing \(\gamma_C\) to a trace;
\item the no-trace/unimodular closure package;
\item a fully explicit statement of what is imported from Paper~6 and what remains open for finite-continuability translation.
\end{enumerate}
This is enough to make Paper 9 a genuine bridge paper.  It would show that the cost language is not merely interpretive: it produces the correct clock sector, identifies the exact spatial obstruction, and supplies a mathematically sharp target for recovering the weak-field two-sector geometry.

\begin{audit}[Do not dilute the result]
The paper should avoid adding speculative cosmology, black holes, or quantum-gravity claims before the scalar trace gate is closed.  The strongest present result is narrow: APF reduces leading weak-field spatial recovery to the no-trace response of a local clock-ruler comparison plane.  That is the clean claim to protect.
\end{audit}

\section{Research tasks for the next draft}

\begin{enumerate}
\item \textbf{Replace the scalar load by a tensorial comparison kernel.}  Define \(Q_{ab}(x)\) as the local cost kernel for infinitesimal comparison continuations and determine how recruitment load perturbs \(Q_{ab}\).
\item \textbf{Derive the response law.}  Start with a nonlinear recruitment functional \(\Erec[Q,\phi]\), expand around a loaded equilibrium, and compute \(\delta^2\Erec/\delta\phi^2\).  The goal is to obtain \(K_t=K_{t0}(1-\alpha\chi)+O(\chi^2)\) rather than assume it.
\item \textbf{Derive the exterior kernel.}  Show that the geometric substrate's unloaded kernel has the Green's-function structure that gives \(\chi\sim1/r\) for a point source in three spatial dimensions.
\item \textbf{Unify signature and dimension.}  Combine finite coordination, irreversible record order, and substrate mode-counting into a single Lorentzian \((1+3)\)-gate.
\item \textbf{Prove the unimodular response lemma.}  Derive \Cref{gate:unimodular-load-response} from the geometric recruitment functional.  By \Cref{prop:unimodular-implies-gamma,cor:spatial-target-equivalence}, this closes \(\gamma_C=1\).
\item \textbf{Derive the gauge-reduced scalar normal form.}  Show from staticity, spherical symmetry, parity, and APF comparison equivalence that all physical first-order scalar response is captured by \(L_{\rm sc}=\lambda_t P_t+\lambda_s P_s\) after coordinate gauge and shear are removed.
\item \textbf{Rule out trace generators.}  Show that static scalar recruitment cannot generate \(\operatorname{tr}L_e\neq0\) unless an independent scalar comparison reservoir is added.  If such a source exists, compute the obstruction coefficient using \Cref{prop:generator-trace-obstruction,prop:trace-obstruction-gamma,prop:normal-form-trace-obstruction}.
\item \textbf{Translate the Paper~6 A9 closure.}  Identify the smallest finite-continuability certificates that imply H1--H3/A9 in the represented smooth regime.
\end{enumerate}

\appendix

\section{Source-of-record redshift calculation}

This appendix records the compact calculation that motivates the first recovery theorem.  Let a static point mass carry an effective accumulated recruitment load \(\chi_{\rm rec}(r)\).  Suppose the exterior linear-response law is
\[
        K_t(r)=K_{t0}(1-\alpha\chi_{\rm rec}(r)),
\]
and suppose the point-source equilibrium plus physical calibration gives
\[
        \alpha\chi_{\rm rec}(r)=\frac{2GM}{rc^2}.
\]
Then
\[
        K_t(r)=K_{t0}\left(1-\frac{2GM}{rc^2}\right).
\]
Using the represented proper-time dictionary,
\[
        d\tau(r)=\sqrt{K_t(r)/K_{t0}}\,dt,
\]
one obtains
\[
        d\tau(r)=\sqrt{1-\frac{2GM}{rc^2}}\,dt.
\]
For two static observers at \(r_B<r_A\), equal coordinate periods correspond to different local proper periods, hence
\[
        \frac{\omega_A}{\omega_B}
        =\frac{d\tau(r_B)}{d\tau(r_A)}
        =\sqrt{\frac{1-2GM/r_Bc^2}{1-2GM/r_Ac^2}}.
\]
Expanding to first order gives
\[
        \frac{\omega_A}{\omega_B}
        \approx
        1-\frac{GM}{c^2}\left(\frac1{r_B}-\frac1{r_A}\right).
\]
This is the static gravitational redshift form.  The calculation is a conditional recovery because the response law, calibration, and proper-time dictionary remain open gates.

\section{Minimal next-lemma inventory}

The next version should try to prove, weaken, or abandon the following lemmas.

\begin{enumerate}
\item \textbf{Kernel lemma.}  Finite local comparison plus isotropic exterior refinement implies a quadratic elliptic kernel with \(1/r\) Green's function in three spatial dimensions.
\item \textbf{Response lemma.}  Loaded-equilibrium expansion of the geometric recruitment functional gives negative linear residual-stiffness response for additional time comparisons.
\item \textbf{Universality lemma.}  All admissible matter clocks couple to the same comparison kernel because composition-dependent coupling would make endpoint comparison non-transitive across sectors.
\item \textbf{Fixed-cell response lemma.}  A static scalar geometric load sources only the trace-free exchange mode \(v_e\), not the trace mode \(u_e\), of the local clock-ruler comparison plane.  This implies \(\mathcal A_e=1+O(c^{-4})\), forcing \(\Theta_s\Theta_t=1+O(c^{-4})\) and hence \(\gamma_C=\gamma_{\rm PPN}=1\).
\item \textbf{Trace-generator lemma.}  Compute the local response generator \(L_e\).  If \(\operatorname{tr}L_e\neq0\), identify \(\zeta_C=\operatorname{tr}L_e\) and test the shifted prediction \(\gamma_C=1-\zeta_C\).
\item \textbf{Self-coupling lemma.}  The cost of geometric distortion itself appears as additional recruitment load, forcing nonlinear closure.
\item \textbf{Ledger-divergence lemma.}  Coordinate relabeling of comparison sites induces an identically conserved geometric ledger, matching stress-energy continuation.
\end{enumerate}

\section{Conclusion}

Paper 10 made the conceptual move: geometry is how finite continuability is carried when the continuation in question is comparison, localization, or transport.  This draft begins the mathematical move: comparison cost first becomes a local cost germ; under convexity and path-additivity gates it admits a Finsler-type length; under the parallelogram certificate it becomes a quadratic metric kernel; committed recruitment load changes residual cost-curvature through an explicit load-probe coupling target; calibrated point-mass load recovers the Schwarzschild time factor and gravitational redshift.  The scalar package also marks the boundary more sharply: clock-sector observables do not determine the spatial response, reciprocal clock-ruler response would, reciprocal response follows from directionwise clock-ruler comparison-cell conservation, cell conservation follows from a fixed-cell variational response principle if scalar recruitment sources only trace-free exchange, all of these are equivalent at first post-Newtonian order to a traceless \(GL(2)\) load-response generator, and the gauge-reduced static scalar normal form shows that, after clock calibration, this trace is the only remaining first-order scalar ambiguity.  If the substrate principles rule out a first-order trace source \(\chi u_e\), then the leading spatial coefficient closes as \(\gamma_C=1\); if they permit such a trace source, \Cref{prop:trace-obstruction-gamma} gives the shifted coefficient \(\gamma_C=1-\zeta_C\).  The next decisive task is therefore no longer to assume the PPN coefficient, but to prove or abandon the unimodular scalar load-response gate.

The result remains intentionally disciplined.  It is not a second derivation of general relativity; Paper~6 already supplies the represented-geometry closure under H1--H3/A9/Lovelock.  Paper~9 is the finite-continuability translation layer: it explains the clock-sector recovery, proves that redshift alone is insufficient, and isolates the no-trace scalar gate whose closure is the local weak-field face of the Paper~6 GR result.

\begin{center}
\fbox{\begin{minipage}{0.88\textwidth}
\centering
\vspace{3pt}
\textbf{Compressed result.}\\[3pt]
Comparison continuations define cost geometry.  Recruitment load perturbs comparison cost.  Calibrated clock-sector response reproduces gravitational time dilation and redshift.  The weak-field spatial sector closes if scalar load transport is unimodular/fixed-cell/trace-free; in the gauge-reduced scalar normal form, the generator trace is the sole PPN-\(\gamma\) obstruction after redshift fixes the clock coefficient.  Paper~6 supplies the represented-GR closure; Paper~9 identifies the finite-continuability weak-field trace gate that must underwrite it from below.\\[3pt]
\end{minipage}}
\end{center}


\paragraph{Current endpoint.}
The paper's present endpoint is therefore precise.  Redshift shows the cost-curvature dictionary has physical bite.  The scalar normal form shows why redshift is not enough.  The no-trace/unimodular gate states the next required theorem in a form that can be attacked mathematically and tested observationally.  Closing that gate is the shortest path from Paper 9 as a recovery program to Paper 9 as a weak-field gravity theorem.

\end{document}
